% coord_xt.tex
%
% Copyright (C) 2022--2026 José A. Navarro Ramón <janr.devel@gmail.com>
% Licencia del código GPLv2
% Licencia Creative Commons Recognition Non-Commercial Share-alike.
% (CC-BY-NC-SA)

\chapter{Otros sistemas 2D: posición, tiempo}
\section{Relatividad de Galileo}
En esta sección vamos a discutir un sistema de coordenadas oblicuo un tanto peculiar. De
hecho, el espacio no será $\symbb{R}^2$, como hace referencia el título del capítulo.
Aun así, será un espacio bidimensional: espacio y tiempo.
En particular tendremos una dimensión de espacio $\symbb{R}$, (una coordenada $x$), y el
tiempo $t$. Tendremos un sistema con dos ejes $x$ y $t$ y lo que mediremos no serán
posiciones o vectores en $\symbb{R}^2$, sino \emph{sucesos}.

\section{Principio de relatividad de Galileo}
El principio de relatividad de Galileo dice que
\begin{quote}
  \emph{%
    Las leyes de la mecánica son idénticas en todos los sistemas de referencia inerciales
    \footnotemark{}
  }.
\end{quote}
Esto significa que ningún experimento realizado en un sistema aislado puede determinar si
está quieto o moviéndose con velocidad constante.
\footnotetext{Marcos de referencia sin aceleración, es decir, con velocidad constante
  (incluido el reposo, $v=0$, que también es una velocidad constante).}

Este principio parte de unas suposiciones
\begin{enumerate}
\item El espacio es absoluto y euclídeo, independiente de los observadores.
\item El tiempo es absoluto. El tiempo es el mismo para todos los observadores,
  independientemente de su movimiento.
\item Existencia de los sistemas de referencia inerciales.
\item La masa de un objeto es la misma para todos los observadores inerciales.
\item Transformación lineal de velocidades. Las velocidades se suman o restan de forma
  lineal. Por ejemplo, si se lanza una piedra a \qty{5}{\kilo\metre\per\hour} desde un
  vehículo a \qty{100}{\kilo\metre\per\hour} en el sentido de avance de este, un
  observador en la carretera medirá la velocidad de la piedra con el valor de
  $5 + 100 = \qty{105}{\kilo\metre\per\hour}$. En particular, la velocidad de la luz
  también está sujeta a esta suposición.
\item Las interacciones entre objetos tiene lugar instantáneamente.
\end{enumerate}

\section{Transformaciones de Galileo}
Las transformaciones de Galileo son un cambio de coordenadas y velocidades que deja
invariante las ecuaciones de Newton.

Como dijimos, nuestro espacio no será $\symbb{R}^2$, sino el de \emph{sucesos}, posición
y tiempo; así, nos centraremos en un eje temporal $t$ y otro espacial $x$, prescindiendo
de las coordenadas $y$ y $z$ para tener una transformación de solo dos variables.

Para centrar ideas, supongamos que tenemos dos observadores, uno $O$ en reposo y el otro
$O'$ que lleva una velocidad constante $u$ con respecto al anterior. En el instante
$t=0$, ambos se encuentran en la misma posición.
El principio de Galileo permite relacionar las mediciones de posición y tiempo de un
suceso, realizadas por ambos observadores $x$, $t$ y $x'$ y $t'$. La posición la
representaremos en el eje de abcisas y el tiempo en ordenadas, como se hace en
relatividad especial y general.

Las transformaciones de Galileo en una dimensión espacial $x$ son
\begin{subequations}
  \begin{align}
    \label{eq:R2-transf_Galileo_x}
    x' &= x - ut\\  
    \label{eq:R2-transf_Galileo_t}
    t' &= t
\end{align}
\end{subequations}
La segunda indica que el tiempo es absoluto.

\subsubsection{Observador $O$ en reposo}
Este observador utiliza un sistema de referencia cartesiano rectangular, ejes $x$ y $t$.
Con él, $O$ puede registrar las posiciones y tiempos del segundo observador o de
cualquier otro suceso.
Como la velocidad de $O'$ es constante, la gráfica del movimiento de $O'$ respecto de $O$
es una línea de pendiente constante. En la figura \ref{fig:R2-oblic_OO'_u} y en las
siguientes, he tomado dos decisiones, teniendo en cuenta que siempre hay sus pros y sus
contras:
\begin{enumerate}
\item La primera es, ¿cuál de los ejes utilizo como abcisas y cual como ordenadas?
  Usaremos como eje de abcisas el eje $x$ y el de ordenadas $t$, al contrario de lo que
  se suele hacer cuando se estudia el movimiento unidimensional en el tiempo en física
  elemental. La explicación es que, como estamos en la \emph{relatividad de Galileo},
  quiero utilizar el esquema que se utiliza en \emph{relatividad especial}, aunque en
  todo lo que implica esta.
  
\item La otra es, ¿qué angulo utilizar en las expresiones matemáticas que obtendremos?
  Teniendo en cuenta que la línea del movimiento relativo de $O'$ respecto de $O$,
  se convertirá en el eje de tiempo $\tilde{t}$ de $O'$ desde el punto de vista de $O$,
  tenemos dos opciones:
  \begin{enumerate}
  \item Utilizar el ángulo $\alpha$ que forma el eje $x$ (abcisas) con la línea de
    pendiente constante. Como hemos dicho, este ángulo nos indicará cómo ve el
    observador $O$, en reposo, los ejes $\tilde{t}$ y $\tilde{x}$ que utiliza $O'$ como
    referencia. El observador en reposo siente que este último no utiliza ejes
    ortogonales, sino que forman un ángulo distinto $\alpha$, en este caso menor de
    \qty{90}{\degree}, porque $O'$ se aleja de $O$ en el sentido positivo del eje $x$
    (sería mayor de \qty{90}{\degree} si se alejara en sentido contrario.
  \item Utilizar el ángulo complementario $\theta = 90-\alpha\,\unit{\degree}$, que
    forma la línea del movimiento de $O'$ con respecto del eje $t$. La pendiente de esta
    línea sería la tangente del ángulo y seo como cuánto se inclina el eje $\tilde{t}$
    de $O'$ con respecto al eje $t$ de $O$. El problema que le veo es que el ángulo, así
    interpretado, debería ser negativo, pues gira en sentido horario. Entiendo que es
    posible no interpretarlo exactamente así y ponerlo positivo, pero no lo considero en
    esta recopilación de notas.
  \end{enumerate}
\end{enumerate}

La figura muestra la velocidad relativa de $O'$ respecto de $O$ es la cotangente de
$\alpha$
\begin{equation}
  \label{eq:R2-transfGalileo_u_alpha}
  u= \dfrac{x}{t} =\cot\alpha
\end{equation}
En esta gráfica, cuanto menor sea $\alpha$, mayor será la velocidad relativa de
$O'$\footnotemark{}.
\footnotetext{Si hubiéramos representado el tiempo en abcisas y la posición en
  ordenadas, $\alpha$ sería la pendiente de la recta, $\alpha=\arctan(x/t)$ y a mayor
  velocidad le correspondería una pendiente mayor.}

Ahora, centramos nuestra atención en un suceso $S$ que tiene lugar en un cierto tiempo $t=t'$, y que es observado tanto por $O$, como por $O'$, pero desde el punto de vista
de $O$, ver figura \ref{fig:R2-oblic_suceso_transf2}.
En ella se comprueba que la distancia $x'$ que mide $O'$ está de acuerdo con la transformación de Galileo \eqref{eq:R2-transf_Galileo_x}.
\begin{figure}[ht]
  \begin{minipage}{0.45\linewidth}
    % Escala
    \def\scl{1.00}
    % 
    \pgfmathsetmacro{\EJESCARTLONG}{3.0}
    % Eje x
    \pgfmathsetmacro{\XLONG}{\EJESCARTLONG}
    % Eje y
    \pgfmathsetmacro{\TLONG}{\EJESCARTLONG}
    % Trayectoria
    \pgfmathsetmacro{\TRAYLONG}{1.0*\EJESCARTLONG}
    \pgfmathsetmacro{\TRAYANG}{60}
    % Punto
    \pgfmathsetmacro{\PMOD}{0.85*\TRAYLONG}
    \pgfmathsetmacro{\PANG}{\TRAYANG}
    % Fondo
    \pgfmathsetmacro{\HORZ}{0.35}
    \pgfmathsetmacro{\VERT}{0.35}
    % 
    \centering
    \begin{tikzpicture}[%
      scale=\scl,
      every node/.style={black,font=\small},
      ejecart/.style={->},
      ejeoblic/.style={->, black},
      proyeccion/.style={black!20, line width=0.4},
      % longitud/.style={line width=.5pt,\colorV},
      texto/.style={black},
      textoapagado/.style={\colorTorig},
      % vtilde/.style={%
      % -{Latex[round,width=4pt,length=6pt]}, line width=1pt, \colorVred},
      pcirculo/.style={fill=\colorPorig, draw=black},
      angulo/.style={fill=green!15, draw=\colorGirodos},
      background/.style={%
        line width=\bgborderwidth,
        draw=\bgbordercolor,
        fill=\bgcolor,
      },
      ]
      % COORDENADAS
      % Origen
      \coordinate (O) at (0,0);
      % Eje x
      \path[save path=\ejex] (O) -- (\XLONG cm, 0) coordinate (xfin);
      % Eje t
      \path[save path=\ejet] (O) -- (0, \TLONG cm) coordinate (tfin);
      % Trayectoria
      \path[save path=\trayect](O) -- ++(\TRAYANG:\TRAYLONG cm)
      coordinate (trayfin);
      % Punto
      \coordinate (P) at (\PANG:\PMOD cm);
      % Proyección de P en los ejes
      % Con eje x
      \path (P) |-  coordinate (px) (xfin);
      \path[save path=\proyx] (P) -- coordinate (pttxt) (px);
      % Posición de texto en eje x
      \path (O) -- coordinate[pos=0.6] (pxtxt) (px);
      % Con eje t
      \path (P) -| coordinate[pos=0.5] (py) (tfin);
      \path[save path=\proyt] (P) -- (py);
      % Posición de texto en eje t
      \path (O) -- coordinate[pos=0.5] (pymed) (py);
      
      % DIBUJOS
      % Ángulo theta
      \path (xfin) -- (O) -- (ytildefin) pic
      [angulo, "\footnotesize $\alpha$",angle radius=11pt, angle eccentricity=0.6]
      {angle = xfin--O--trayfin};
      % Proyecciones de componentes
      % Eje x
      \draw[proyeccion, use path=\proyx];
      \node[below] at (pxtxt) {$x=ut$};
      % Eje t
      \draw[proyeccion, use path=\proyt];
      \node[right] at (pttxt) {$t$};

      % Ejes
      % x
      \draw[ejecart, use path=\ejex];
      \node[below=1.2pt] at (xfin) {\small $x$};
      % t
      \draw[ejecart, use path=\ejey];
      \node[above] at (yfin) {$t$};
      % Trayectoria
      \draw[line width=1pt, use path=\trayect];
      
      % Texto informativo
      \node[above right=0pt and -1pt] at (trayfin)
      {$\left[\cot\alpha=\dfrac{x}{t} = \dfrac{ut}{t} = u\right]$};
      % Punto
      \fill (P) circle[radius=1pt];
      % Origen
      \filldraw (O) circle [radius=.2pt];
      
      % Fondo amarillo
      \coordinate (SW) at ($(current bounding box.south west) + (-\HORZ cm,-\VERT cm)$);
      \coordinate (NE) at ($(current bounding box.north east) + (\HORZ cm,\VERT cm)$);
      \begin{scope}[on background layer]
        \draw[background] (SW) rectangle (NE);
      \end{scope}
    \end{tikzpicture}
    \caption{Gráfica del movimiento uniforme de $O'$ con respecto de $O$. Es una línea
      recta de pendiente $\tan\alpha = t/x$ o $\cot\alpha = x/t$.
      La velocidad constante de $O'$ con respecto de $O$ es $u=\cot\alpha$.
      La distancia $x$ representa la distancia entre $O$ y $O'$ en cada instante $t$.}
    \label{fig:R2-oblic_OO'_u}
  \end{minipage}
  \hspace{1em}
  \begin{minipage}{0.45\linewidth}
    % Escala
    \def\scl{1.00}
    % 
    \pgfmathsetmacro{\EJESCARTLONG}{3.0}
    % Eje x
    \pgfmathsetmacro{\XLONG}{\EJESCARTLONG}
    % Eje y
    \pgfmathsetmacro{\TLONG}{\EJESCARTLONG}
    % Trayectoria
    \pgfmathsetmacro{\TRAYLONG}{1.0*\EJESCARTLONG}
    \pgfmathsetmacro{\TRAYANG}{60}
    % Punto
    \pgfmathsetmacro{\PMOD}{0.80*\TRAYLONG}
    \pgfmathsetmacro{\PANG}{\TRAYANG}
    % Suceso
    \pgfmathsetmacro{\OPRIMASUCESOLONG}{0.60*(\XLONG - sin(\PANG))}
    \pgfmathsetmacro{\BELOWSUCESO}{0.08}
    % Fondo
    \pgfmathsetmacro{\HORZ}{0.50}
    \pgfmathsetmacro{\VERT}{0.35}
    % 
    \centering
    \begin{tikzpicture}[%
      scale=\scl,
      every node/.style={black,font=\small},
      ejecart/.style={->},
      ejeoblic/.style={->, black},
      proyeccion/.style={black!20, line width=0.4pt},
      trayectoria/.style={black!25, line width=1pt},
      % longitud/.style={line width=.5pt,\colorV},
      puntosuceso/.style={fill=red, draw=black},
      OOprima/.style={black!80, line width=1pt},
      OprimaSuceso/.style={red, line width=1pt},
      OSuceso/.style={green!70!black, line width=1pt},
      texto/.style={black},
      textoapagado/.style={\colorTorig},
      % vtilde/.style={%
      % -{Latex[round,width=4pt,length=6pt]}, line width=1pt, \colorVred},
      pcirculo/.style={fill=\colorPorig, draw=black},
      angulo/.style={fill=green!15, draw=\colorGirodos},
      background/.style={%
        line width=\bgborderwidth,
        draw=\bgbordercolor,
        fill=\bgcolor,
      },
      ]
      % COORDENADAS
      % Origen
      \coordinate (O) at (0,0);
      % Eje x
      \path[save path=\ejex] (O) -- (\XLONG cm, 0) coordinate (xfin);
      % Eje t
      \path[save path=\ejet] (O) -- (0, \TLONG cm) coordinate (tfin);
      % Trayectoria
      \path[save path=\trayect](O) -- ++(\TRAYANG:\TRAYLONG cm)
      coordinate (trayfin);
      % Punto
      \coordinate (P) at (\PANG:\PMOD cm);
      % Proyección de P en los ejes
      % Con eje x
      \path (P) |-  coordinate (px) (xfin);
      \path[save path=\proyx] (P) -- coordinate (pttxt) (px);
      % Posición de texto en eje x
      \path (O) -- coordinate[pos=0.6] (pxtxt) (px);
      % Con eje t
      \path (P) -| coordinate[pos=0.5] (py) (tfin);
      \path[save path=\OOprima] (P) -- coordinate[pos=0.5] (OOprimatxt) (py);
      % Posición de texto en eje t
      \path (O) -- coordinate[pos=0.5] (pymed) (py);
      % Posición del suceso (S)
      \path (P) -- ++(\OPRIMASUCESOLONG, 0) coordinate (S);
      % Distancia O' Suceso
      \path[save path=\OprimaSuceso]
      (P) -- coordinate[pos=0.5] (OprimaStxt) ++(\OPRIMASUCESOLONG, 0);
      % Posición bajo el suceso (belowS)
      \path (S) -- ++(0, -\BELOWSUCESO) coordinate (belowS);
      \path (belowS) -| coordinate[pos=0.5] (belowO) (tfin);
      % Distancia belowO belowSuceso
      \path[save path=\belowOSuceso] (belowO) -- coordinate[pos=0.5] (OStxt) (belowS);
      
      % DIBUJOS
      % Ángulo theta
      \path (xfin) -- (O) -- (ytildefin) pic
      [angulo, "\footnotesize $\alpha$",angle radius=11pt, angle eccentricity=0.6]
      {angle = xfin--O--trayfin};
      
      % Distancia OO'
      \draw[OOprima, use path=\OOprima];
      \node[above] at (OOprimatxt) {$ut$};
      % Distancia O'S
      \draw[OprimaSuceso, use path=\OprimaSuceso];
      \node[above, red!70!black] at (OprimaStxt) {$x'$};
      
      % Ejes
      % x
      \draw[ejecart, use path=\ejex];
      \node[below=1.2pt] at (xfin) {\small $x$};
      % t
      \draw[ejecart, use path=\ejey];
      \node[above] at (yfin) {$t$};
      % Trayectoria
      \draw[trayectoria, use path=\trayect];

      % Distancia OS
      \draw[OSuceso, use path=\belowOSuceso]; 
      \node[below, green!60!black] at (OStxt) {$x$};

      % Texto informativo
      \node[above right=20pt and 36pt] at (P)
      {$\left[
          \color{red!80!black}{x'}
          \color{black}{=}
          \color{green!60!black}{x}
          \color{black}{-ut}
        \right]$};
      
      % Suceso
      \filldraw[puntosuceso] (S) circle[radius=1pt];
      \node[above right] at (S) {\footnotesize $S$};
      % Punto
      \fill (P) circle[radius=1pt];
      % Origen
      \filldraw (O) circle [radius=.2pt];
      
      % Fondo amarillo
      \coordinate (SW) at ($(current bounding box.south west) + (-\HORZ cm,-\VERT cm)$);
      \coordinate (NE) at ($(current bounding box.north east) + (\HORZ cm,\VERT cm)$);
      \begin{scope}[on background layer]
        \draw[background] (SW) rectangle (NE);
      \end{scope}
    \end{tikzpicture}
    \caption{Comprobación de una de las transformaciones de Galileo, la que relaciona
      la posición de un cierto suceso medido por ambos, $O$ y $O'$ en un instante $t=t'$.
      El tiempo absoluto está implícito en la esta figura.}
    \label{fig:R2-oblic_suceso_transf2}
  \end{minipage}
\end{figure}

Las transformaciones de Galileo en forma matricial son\footnotemark{}
\footnotetext{En relatividad especial (Einstein), es costumbre cambiar el orden y
  escribir primero la componente temporal y después las espaciales, dando a la primera
  el índice cero. Aquí, seguiremos este criterio, para familiarizarnos con él, no porque
  sea necesario, sino porque es el estándar en relatividad einsteniana.}
\[
  \begin{pmatrix}
    t'\\
    x'
  \end{pmatrix}
  = \begin{pmatrix}
    1 & 0 \\
    -u & 1
  \end{pmatrix}
  \,
  \begin{pmatrix}
    t\\
    x
  \end{pmatrix}
\]

Estas componentes son contravariantes, debido a que si ampliamos la escala de medida de posiciones y/o tiempo, la componente respectiva se hace más pequeña, y viceversa.
Por tanto, podremos escribir $x^0 = t$, $x^1 = x$, $\tilde{x}^0 = t'$ y
$\tilde{x}^1 = x'$.
Las transformaciones de Galileo quedarían
\begin{equation}
  \label{eq:R2-transf_galileo}
  \begin{pmatrix}
    \tilde{x}^0\\
    \tilde{x}^1
  \end{pmatrix}
  = \begin{pmatrix}
    1 & 0 \\
    -u & 1
  \end{pmatrix}
  \,
  \begin{pmatrix}
    x^0\\
    x^1
  \end{pmatrix}
\end{equation}
donde $u$ es la velocidad relativa de $O'$ con respecto de $O$.

De \eqref{eq:R2-transf_galileo}, podemos deducir que la matriz de transformación es la
inversa del cambio de base $\mmm{M}^{-1}$
\begin{equation}
  \mmm{M}^{-1}
  =
  \begin{pmatrix}
    1 & 0\\
    -u & 1
  \end{pmatrix}
  =
  \begin{pmatrix}
    1 & 0\\
    -\cot\alpha & 1
  \end{pmatrix}
\end{equation}

La matriz inversa de la anterior es la matriz de cambio de base
\begin{equation}
  \label{eq:R2-M}
  \mmm{M}
  =
  \begin{pmatrix}
    1 & 0\\
    u & 1
  \end{pmatrix}
  =
  \begin{pmatrix}
    1 & 0\\
    \cot\alpha & 1
  \end{pmatrix}
\end{equation}

\section{Bases y matriz de cambio de base}
El objetivo de este apartado es el de obtener, de una forma equivalente, la matriz de
cambio de base, hallada en \eqref{eq:R2-M}, partiendo de las transformaciones de Galileo.

\subsubsection{Base $B$ para el observador $O$ en reposo}
Desde el punto de vista de $O$, que se considera en reposo, utiliza un sistema de
coordenadas rectangular con base
\begin{equation}
  B
  = \set{\xhat{e}_0, \xhat{e}_1}
  = \set{\xhat{e}_t, \xhat{e}_x}
  = \set{\xhat{e}_t, \uvec{\i}}
\end{equation}
donde el vector unitario $\xhat{e}_0$ está en unidades de tiempo, por ejemplo,
\qty{1}{\second}, y el vector unitario $\xhat{e}_1$ está en unidades de longitud, por
ejemplo \qty{1}{\metre}; luego no son equivalentes, como ocurría en $\symbb{R}^2$.
En cualquier caso, prescindiendo de las unidades, el módulo de estos dos vectores vale 1
\begin{subequations}
\begin{align}
  |\xhat{e}_0| &= |\xhat{e}_t| = 1\\
  |\xhat{e}_1| &= |\xhat{e}_x| = |\uvec{\i}| = 1
\end{align}
\end{subequations}

Según lo visto en las figuras anteriores, el observador $O'$, que se considera en
movimiento, desde el punto de vista de $O$ tiene un eje de tiempo $t'$ inclinado un
ángulo $\alpha$ que depende de la velocidad relativa $u$
(ver \eqref{eq:R2-transfGalileo_u_alpha}); en cambio, el eje $x'$ coincide con el de
$O$.

Nótese que, $O'$ tiene derecho, como sistema inercial, a considerarse en reposo, y desde
este punto de vista de $O'$, el eje del tiempo de $O$ se inclinaría en sentido
antihorario formando un ángulo de $\pi/2+\alpha$ con el eje $x$.

En este texto supondremos siempre que $O$ es el que se encuentra en reposo (utilizamos
el punto de vista de $O$).

\subsubsection{Base $B'$ para el observador $O'$ desde el punto de vista de $O$}
La base utilizada por $O'$, sería
\begin{equation}
  B'
  = \set{\xtilde{e}_0, \xtilde{e}_1}
  = \set{\xtilde{e}_{t'}, \xhat{e}_{x'}}
\end{equation}

¿Qué módulo tienen estos vectores de $B'$? El del eje $x'$ es la unidad, igual que el
de la base $B$ de $O$. Pero, debido a que el tiempo es absoluto, $t=t'$, las líneas
paralelas al eje $x$ son líneas de igual tiempo $t=t'$, el vector temporal de la base $B'$ tiene
que ser mayor que el vector de la base temporal de $O$, que vale 1, ver figura
\ref{fig:R2-oblic_etprima_et}
\[
  \cos(\pi/2 - \alpha)
  = \sin\alpha
  = \dfrac{|\xhat{e}_t|}{|\xtilde{e}_{t'}|}
  = \dfrac{1}{|\xtilde{e}_{t'}|}
\]
De donde
\[
  |\xtilde{e}_{t'}| = \dfrac{1}{\sin\alpha} = \csc\alpha
\]

\begin{subequations}
\begin{align}
  |\xtilde{e}_0| &= |\tilde{e}_{t'}| = \csc\alpha = \sqrt{1+u^2}\\
  |\xtilde{e}_1| &= |\xhat{e}_{x'}| = 1
\end{align}
\end{subequations}

\begin{figure}[ht]
  \begin{minipage}{0.45\linewidth}
    % Escala
    \def\scl{1.00}
    % 
    \pgfmathsetmacro{\EJESCARTLONG}{3.0}
    \pgfmathsetmacro{\EJEOBLICLONG}{2.7}
    % Eje x
    \pgfmathsetmacro{\XLONG}{\EJESCARTLONG}
    % Eje y
    \pgfmathsetmacro{\YLONG}{\EJESCARTLONG}
    % Eje xtilde
    \pgfmathsetmacro{\XTILDEANG}{0}
    \pgfmathsetmacro{\XTILDELONG}{\EJEOBLICLONG}    
    % Eje ytilde
    \pgfmathsetmacro{\YTILDEANG}{60}
    \pgfmathsetmacro{\YTILDELONG}{\EJESCARTLONG}
    % Longitud del vector unidad en el eje x': etilde1
    \pgfmathsetmacro{\EXTILDE}{0.70*\EJEOBLICLONG}
    \pgfmathsetmacro{\ETTILDE}{\EXTILDE/(sin(\YTILDEANG))}
    % Vector P
    %\pgfmathsetmacro{\PMOD}{2.3}
    %\pgfmathsetmacro{\PANG}{40}
    % Fondo
    \pgfmathsetmacro{\HORZ}{0.65}
    \pgfmathsetmacro{\VERT}{0.25}
    % 
    \centering
    \begin{tikzpicture}[%
      scale=\scl,
      every node/.style={black,font=\small},
      ejecart/.style={->, black!30},
      ejeoblic/.style={->, black},
      proyeccion/.style={black!20, line width=0.4},
      % longitud/.style={line width=.5pt,\colorV},
      texto/.style={black},
      textoapagado/.style={\colorTorig},
      vector/.style={-{Latex[round]}, shorten >=1.2pt, line width=.8pt,\colorV},
      versor/.style={-{Latex[round, width=4pt, length=6pt]}, line width=1pt, \colorVred},
      % vtilde/.style={%
      % -{Latex[round,width=4pt,length=6pt]}, line width=1pt, \colorVred},
      pcirculo/.style={fill=\colorPorig, draw=black},
      angulo/.style={fill=green!15, draw=\colorGirodos},
      background/.style={%
        line width=\bgborderwidth,
        draw=\bgbordercolor,
        fill=\bgcolor,
      },
      ]
      % COORDENADAS
      % Origen
      \coordinate (O) at (0,0);
      % Eje x
      \path[save path=\ejex] (O) -- (\XLONG cm, 0) coordinate (xfin);
      % Eje y
      \path[save path=\ejey] (O) -- (0, \YLONG cm) coordinate (yfin);
      % Eje xtilde
      \path[save path=\ejextilde] (O) -- ++(\XTILDEANG:\XTILDELONG cm)
      coordinate (xtildefin);
      \node[texto, below] at (xtildefin) {\footnotesize $\tilde{x}$};
      % Eje ytilde
      \path[save path=\ejeytilde,name path=O--ytilde](O) -- ++(\YTILDEANG:\YTILDELONG cm)
      coordinate (ytildefin);
      % Vector
      %\path[save path=\r] (O) -- coordinate[pos=0.5] (pmid) (\PANG:\PMOD cm)
      %coordinate (P);
      % Componentes covariantes y contravariantes
      % Intersección de horizontal desde P hasta ytilde
      % -> Componente contravariante eje ytilde (CTV2) 
      \path[name path=P--y] (P) -| (yfin);
      \path[name intersections={of=P--y and O--ytilde,by=CTV2}];
      % Paralela (CTV2) -- (O) que pase por (P)
      % -> Componente contravariante eje xtilde (CTV1)
      \path (P) -- ($(P) - (CTV2) - (O)$) coordinate (CTV1);
      % Componente covariante con eje xtilde (COV1)
      \path (P) |- coordinate[pos=0.5] (COV1) (xtildefin);
      % Perpendicular a O--ytilde que pasa por P
      % -> Componente covariante con eje ytilde (COV2)
      \path (P) -- ($(O)!(P)!(ytildefin)$) coordinate (COV2);
      % Proyecciones de las componentes
      
      % DIBUJOS
      % Ángulo theta
      \path (xfin) -- (O) -- (ytildefin) pic
      [angulo, "\footnotesize $\alpha$",angle radius=12pt, angle eccentricity=0.65]
      {angle = xfin--O--trayfin};p

      % Ejes
      % x
      \draw[ejecart, use path=\ejex];
      \node[textoapagado, below=1.2pt] at (xfin) {\small $x$};
      % y
      \draw[ejecart, use path=\ejey];
      \node[above, textoapagado] at (yfin) {$y$};
      % xtilde
      \draw[ejeoblic, use path=\ejextilde];
      % \node[below, texto] at (xtildefin) {$\tilde{x}^1$};
      % ytilde
      \draw[ejeoblic, use path=\ejeytilde];
      \node[above right=0pt and -4pt, texto] at (ytildefin) {$\tilde{y}$};

      % Versor eT en gris
      \draw[versor, gray] (O) -- node[gray, pos=0.6, left]
        {$\xhat{e}_t$} ++(0, \EXTILDE cm) coordinate (e0fin);
      % Versor etilde0
      \draw[versor] (O) -- node[\colorTred, pos=0.55, right] {$\xtilde{e}_{t'}$}
      ++(\YTILDEANG:\ETTILDE cm) coordinate (etilde0fin);
      % Versor etilde1
      \draw[versor] (O) -- node[\colorTred, pos=0.6, below]
      {$\xhat{e}_{x'} = \color{black!50}{\uvec{\i}}$}
      ++(\EXTILDE cm, 0) coordinate (etilde1fin);

      % Línea horizontal que une e0 con etilde0
      \draw[proyeccion] (e0fin) -- (etilde0fin);
      % Extensión izda de la línea
      \draw[proyeccion] (e0fin) --  ++(-0.8, 0);
      % Extensión derecha de la línea
      \draw[proyeccion] (etilde0fin) --
      node[pos=0.6, above] {\footnotesize $\color{gray}{t = t' = \qty{1}{\second}}$}
      ++(1.8, 0);
      
      % Origen
      \filldraw (O) circle [radius=.2pt];
      
      % Fondo amarillo
      \coordinate (SW) at ($(current bounding box.south west) + (-\HORZ cm,-\VERT cm)$);
      \coordinate (NE) at ($(current bounding box.north east) + (\HORZ cm,\VERT cm)$);
      \begin{scope}[on background layer]
        \draw[background] (SW) rectangle (NE);
      \end{scope}
    \end{tikzpicture}
    \caption{En rojo, los vectores de la base $B'$, donde $\xtilde{e}_{t'}$ se ve
      inclinado, desde el punto de vista del observador en reposo $O$. En gris, los
      vectores unitario temporal y espacial de la base $B$. Todas las líneas
      horizontales que se puedan trazar son equitemporales. La que se muestra indica
      $t=t'=\qty{1}{\second}$.}
    \label{fig:R2-oblic_etprima_et}
  \end{minipage}
  \hspace{1em}
  \begin{minipage}{0.45\linewidth}
    % Escala
    \def\scl{1.00}
    % 
    \pgfmathsetmacro{\EJESCARTLONG}{3.0}
    \pgfmathsetmacro{\EJEOBLICLONG}{2.7}
    % Eje x
    \pgfmathsetmacro{\XLONG}{\EJESCARTLONG}
    % Eje y
    \pgfmathsetmacro{\YLONG}{\EJESCARTLONG}
    % Eje xtilde
    \pgfmathsetmacro{\XTILDEANG}{0}
    \pgfmathsetmacro{\XTILDELONG}{\EJEOBLICLONG}    
    % Eje ytilde
    \pgfmathsetmacro{\YTILDEANG}{60}
    \pgfmathsetmacro{\YTILDELONG}{\EJESCARTLONG}
    % Longitud del vector unidad en el eje x': etilde1
    \pgfmathsetmacro{\EXTILDE}{0.7*\EJEOBLICLONG}
    \pgfmathsetmacro{\ETTILDE}{\EXTILDE/(sin(\YTILDEANG))}
    % Vector P
    %\pgfmathsetmacro{\PMOD}{2.3}
    %\pgfmathsetmacro{\PANG}{40}
    % Fondo
    \pgfmathsetmacro{\HORZ}{0.80}
    \pgfmathsetmacro{\VERT}{0.25}
    % 
    \centering
    \begin{tikzpicture}[%
      scale=\scl,
      every node/.style={black,font=\small},
      ejecart/.style={->, black!30},
      ejeoblic/.style={->, black},
      proyeccion/.style={black!20, line width=0.4},
      % longitud/.style={line width=.5pt,\colorV},
      texto/.style={black},
      textoapagado/.style={\colorTorig},
      vector/.style={-{Latex[round]}, shorten >=1.2pt, line width=.8pt,\colorV},
      versor/.style={-{Latex[round, width=4pt, length=6pt]}, line width=1pt, \colorVred},
      % vtilde/.style={%
      % -{Latex[round,width=4pt,length=6pt]}, line width=1pt, \colorVred},
      pcirculo/.style={fill=\colorPorig, draw=black},
      angulo/.style={fill=green!15, draw=\colorGirodos},
      background/.style={%
        line width=\bgborderwidth,
        draw=\bgbordercolor,
        fill=\bgcolor,
      },
      ]
      % COORDENADAS
      % Origen
      \coordinate (O) at (0,0);
      % Eje x
      \path[save path=\ejex] (O) -- (\XLONG cm, 0) coordinate (xfin);
      % Eje y
      \path[save path=\ejey] (O) -- (0, \YLONG cm) coordinate (yfin);
      % Eje xtilde
      \path[save path=\ejextilde] (O) -- ++(\XTILDEANG:\XTILDELONG cm)
      coordinate (xtildefin);
      \node[texto, below] at (xtildefin) {\footnotesize $\tilde{x}$};
      % Eje ytilde
      \path[save path=\ejeytilde,name path=O--ytilde](O) -- ++(\YTILDEANG:\YTILDELONG cm)
      coordinate (ytildefin);
      % Componente y
      \path (O) -- coordinate[pos=0.5] (compy) (e0fin);
      
      % DIBUJOS
      % Ángulo theta
      \path (xfin) -- (O) -- (ytildefin) pic
      [angulo, "\footnotesize $\alpha$",angle radius=11pt, angle eccentricity=0.63]
      {angle = xfin--O--trayfin};

      % Línea horizontal que une e0 con etilde0
      \draw[proyeccion] (e0fin) -- (etilde0fin);
      % Línea vertical de proyección de etilde0
      \draw[proyeccion] (etilde0fin) |- coordinate[pos=0.5] (compx) (xtildefin);
      
      % Texto componente x
      \path (O) -- node[pos=0.5, below]
      {\footnotesize $|\xtilde{e}_{t'}| \cos\alpha$} (compx);
      % Texto componente y
      \node[rotate=90, above left=18pt and 0pt] at (compy)
      {\footnotesize $|\xtilde{e}_{t'}| \sin\alpha$};      
      
      % Ejes
      % x
      \draw[ejecart, use path=\ejex];
      \node[textoapagado, below=1.2pt] at (xfin) {\small $x$};
      % y
      \draw[ejecart, use path=\ejey];
      \node[above, textoapagado] at (yfin) {$y$};

      % xtilde
      \draw[ejeoblic, use path=\ejextilde];
      % ytilde
      \draw[ejeoblic, use path=\ejeytilde];
      \node[above right=0pt and -4pt, texto] at (ytildefin) {$\tilde{y}$};

      % Versor etilde0
      \draw[versor] (O) -- node[\colorTred,pos=0.65, left] {\small $\xtilde{e}_{t'}$}
      ++(\YTILDEANG:\ETTILDE cm) coordinate (etilde0fin);
      % Versor etilde1
      %\draw[versor] (O) -- node[\colorTred, pos=0.6, below] {\scriptsize $\xhat{e}_1$}
      %++(\EXTILDE cm, 0) coordinate (etilde1fin);

      % Origen
      \filldraw (O) circle [radius=.2pt];
      
      % Fondo amarillo
      \coordinate (SW) at ($(current bounding box.south west) + (-\HORZ cm,-\VERT cm)$);
      \coordinate (NE) at ($(current bounding box.north east) + (\HORZ cm,\VERT cm)$);
      \begin{scope}[on background layer]
        \draw[background] (SW) rectangle (NE);
      \end{scope}
    \end{tikzpicture}
    \caption{Componente espacial y temporal en la base $B$ del vector temporal de la
      base $B'$, $\xtilde{e}_{t'}$, en rojo.
      Como su módulo es $|\xtilde{e}_{t'}|= \csc\alpha$, se deduce que
      $\xtilde{e}_{t'} = \cot\alpha\,\uvec{i} + \xhat{e}_t$ en la base $B$. Esto nos
      servirá para deducir la matriz de cambio de base desde el punto de vista de $O$.}
    \label{fig:R2-oblic_etprima}
  \end{minipage}
\end{figure}

Resumiendo, los vectores de la base $B'$ en función de $B$ valen
\begin{align}
  \label{eq:R2-oblic_cov_e0_et}
  \xtilde{e}_0
  &=
    \xtilde{e}_{t'}
    = \xhat{e}_t + \cot\alpha\,\uvec{\i}
    = (1, \cot\alpha)
    = \xhat{e}_t + u\,\uvec{\i}
    =(1, u)\\
  \label{eq:R2-oblic_cov_e1_ex}
  \xhat{e}_1
  &=
    \xhat{e}_{x'}
    = 0\xhat{e}_t+\xhat{e}_x
    = 0\xhat{e}_t + \uvec{\i}
    = \uvec{\i}
    = (0, 1)
\end{align}

La matriz de cambio de base se construye situando las componentes de los vectores de la
base $B'$ en función de $B$ en columnas. Así $B' = B \mmm{M}$. A continuación
expresamos lo mismo, tanto en función de $\alpha$, como de $u$, con distinta
nomenclatura de índices
\begin{subequations}
  \begin{align}
    %\label{eq:R2-oblic_cambio_base_01alpha}
    \begin{pmatrix}
      \xtilde{e}_0 & \xtilde{e}_1
    \end{pmatrix}
    &=
    \begin{pmatrix}
      \xhat{e}_0 & \xhat{e}_1
    \end{pmatrix}
    \,
    \begin{pmatrix}
      1 & 0\\
      \cot\alpha & 1
    \end{pmatrix}\\
    %\label{eq:R2-oblic_cambio_base_txu}
    \begin{pmatrix}
      \xtilde{e}_{t'} & \xtilde{e}_{x'}
    \end{pmatrix}
    &=
      \begin{pmatrix}
        \xhat{e}_{t} & \uvec{\i}
      \end{pmatrix}
      \,
      \begin{pmatrix}
        1 & 0\\
        u & 1
      \end{pmatrix}
  \end{align}
\end{subequations}
La matriz de cambio de base coincide con la encontrada antes \eqref{eq:R2-M}.

\section{Métrica en $O'$ desde el punto de vista de $O$}
Calcularemos la métrica en a partir del producto escalar de los vectores de $B'$,
y teniendo en cuenta que los vectores $\xhat{e}_t$ y $\uvec{\i}$ son ortogonales y
unitarios
{\small
\begin{subequations}
  \begin{align}
    \xtilde{e}_0\cdot\xtilde{e}_0
    &= \xtilde{e}_{t'}\cdot\xtilde{e}_{t'}
      = (\xhat{e}_t+\cot\alpha\,\uvec{\i})\cdot(\xhat{e}_t+\cot\alpha\,\uvec{\i})
      = 1 + \cot^2\alpha = \csc^2\alpha\\
    \xtilde{e}_0\cdot\xtilde{e}_1
    &= \xtilde{e}_{t'}\cdot\xtilde{e}_{x'}
      = (\xhat{e}_t + \cot\alpha\,\uvec{\i})\cdot\uvec{\i}=\cot\alpha\\
    \xtilde{e}_1\cdot\xtilde{e}_0
    &= \xtilde{e}_{x'}\cdot\xtilde{e}_{t'}= \cdots = \cot\alpha\\
    \xtilde{e}_1\cdot\xtilde{e}_1
    &= \xtilde{e}_{x'}\cdot\xtilde{e}_{x'}
      = \uvec{\i}\cdot\uvec{\i}
      = 1
  \end{align}
\end{subequations}
}

Como resultado, el tensor métrico en $O'$ desde el punto de vista de $O$ se puede
representar como
{\small
  \[
    \mmm{g}_{ij}
    = \begin{pmatrix}
      g_{00} & g_{01}\\
      g_{10} & g_{11}
    \end{pmatrix}
    = \begin{pmatrix}
      \xtilde{e}_0\cdot\xtilde{e}_0 & \xtilde{e}_0\cdot\xtilde{e}_1\\
      \xtilde{e}_1\cdot\xtilde{e}_0 & \xtilde{e}_1\cdot\xtilde{e}_1
    \end{pmatrix}
    = \begin{pmatrix}
      \xtilde{e}_{t'}\cdot\xtilde{e}_{t'} & \xtilde{e}_{t'}\cdot\xtilde{e}_{x'}\\
      \xtilde{e}_{x'}\cdot\xtilde{e}_{t'} & \xtilde{e}_{x'}\cdot\xtilde{e}_{x'}
    \end{pmatrix}
    = \begin{pmatrix}
      \csc^2\alpha & \cot\alpha\\
      \cot\alpha & 1
    \end{pmatrix}
  \]
}

Como $u = \cot\alpha$ y $\csc^2\alpha$, según \ref{fig:R2-oblic_OO'_u}
\begin{align*}
  \csc\alpha
  &=
    \dfrac{\sqrt{t^2+u^2t^2}}{t}
    = \dfrac{\sqrt{t^2(1+u^2)}}{t}
    = \dfrac{\cancelout{t}\sqrt{1+u^2}}{\cancelout{t}}
    = \sqrt{1+u^2}
\end{align*}

Entonces, la métrica se puede escribir como
\begin{equation}
  \mmm{g}_{ij}
  = \begin{pmatrix}
    \csc^2\alpha & \cot\alpha\\
    \cot\alpha & 1      
  \end{pmatrix}
  = \begin{pmatrix}
    1+u^2 & u\\
    u & 1      
  \end{pmatrix}
\end{equation}
Su inversa sería
{\small
  \begin{align*}
    \mmm{g}^{ij}
    &=
      \dfrac{1}{\csc^2\alpha-\cot^2\alpha}
      \begin{pmatrix}
        1 & -\cot\alpha\\
        -\cot\alpha & \csc^2\alpha      
      \end{pmatrix}
      = \dfrac{1}{1}
      \begin{pmatrix}
        1 & -\cot\alpha\\
        -\cot\alpha & \csc^2\alpha
      \end{pmatrix}
      = \begin{pmatrix}
        1 & -\cot\alpha\\
        -\cot\alpha & \csc^2\alpha
      \end{pmatrix}    
  \end{align*}
}
Resultando
\begin{equation}
  \mmm{g}^{ij}
    = \begin{pmatrix}
      1 & -\cot\alpha\\
      -\cot\alpha & \csc^2\alpha
    \end{pmatrix}
    = \begin{pmatrix}
      1 & -u\\
      -u & 1+u^2
    \end{pmatrix}
\end{equation}

\section{Vectores covariantes y contravariantes}
Hasta ahora, hemos tratado con los vectores covariantes y las componentes contravariantes
en $O'$ desde el punto de vista de $O$.
En notación de Einstein
\[
  \vvv{v} = x^i\,\tilde{e}_i
  \hspace{0.75em}\text{con } i = 0, 1
\]

Vamos a subir los índices de los vectores de $B$, $\xtilde{e}_i = g^{ij}\xtilde{e}_j$
En la base $B'$
\begin{align*}
  \xtilde{e}^0
  &=
    g^{00}\xtilde{e}_0 + g^{01}\xtilde{e}_1
    = 1\xtilde{e}_0 - u\xtilde{e}_1
    = \xtilde{e}_0 - u\xtilde{e}_1\\
  \xtilde{e}^1
  &=
    g^{10}\xtilde{e}_0 + g^{11}\xtilde{e}_1
    = -u\xtilde{e}_0 + (1+u^2)\xtilde{e}_1
\end{align*}

En la base $B$, teniendo en cuenta \eqref{eq:R2-oblic_cov_e0_et}
\eqref{eq:R2-oblic_cov_e1_ex}
\begin{align}
  \xtilde{e}^0
  &=
    \xtilde{e}_0 - u\xtilde{e}_1
    = (\xhat{e}_t + u\,\uvec{\i}) - u\,\uvec{\i}
    = \xhat{e}_t\\
  \xtilde{e}^1
  &=
    -u\xtilde{e}_0 + (1+u^2)\xtilde{e}_1
    -u (\xhat{e}_t + u\,\uvec{\i}) + (1 + u^2)\,\uvec{\i}
    = -u\xhat{e}_t + \uvec{\i}
\end{align}
y volvemos a escribir los resultados para poder compararlas en la misma página, todas
con respecto a la base rectangular $B$.

Vectores covariantes
\begin{subequations}
\begin{align}
  \xtilde{e}_0 &= \xtilde{e}_{t'} = \xhat{e}_{t} + u\,\uvec{\i}\\
  \xtilde{e}_1 &= \xtilde{e}_{x'} = \uvec{\i}
\end{align}
\end{subequations}

Vectores contravariantes
\begin{subequations}
\begin{align}
  \xtilde{e}^0 &= \xtilde{e}^{t'} = \xhat{e}_t\\
  \xtilde{e}^1 &= \xtilde{e}^{x'} = -u\xhat{e}_t + \uvec{\i}
\end{align}
\end{subequations}


%
%\section{Coordenadas polares planas}
%Aquí se ponen las cosas algo más interesantes. En este sistema, se utilizan dos
%coordenadas, $r$ y $\theta$. La primera coordenada es la distancia desde el punto del espacio al origen de coordenadas, y la segunda es el ángulo que forma el eje de abcisas
%con la línea que une el origen de coordenadas con el punto, ver figura
%\ref{fig:R2-coord_polares}.
%
%\begin{figure}[ht]
%  \centering
%  \begin{minipage}{0.40\linewidth}
%    % Escala
%    \def\scl{1.25}
%    % 
%    \pgfmathsetmacro{\EJESEXTRA}{-0.2}
%    \pgfmathsetmacro{\EJESLONG}{2.3}
%    % Eje x
%    \pgfmathsetmacro{\XMLONG}{\EJESEXTRA}
%    \pgfmathsetmacro{\XPLONG}{\EJESLONG}
%    % Eje y
%    \pgfmathsetmacro{\YMLONG}{\EJESEXTRA}
%    \pgfmathsetmacro{\YPLONG}{\EJESLONG}
%    % Vector P
%    \pgfmathsetmacro{\PMOD}{2.5}
%    \pgfmathsetmacro{\PANG}{40}
%    % Fondo
%    \pgfmathsetmacro{\HORZ}{0.25}
%    \pgfmathsetmacro{\VERT}{0.25}
%    % 
%    \centering
%    \begin{tikzpicture}[%
%      scale=\scl,
%      every node/.style={black,font=\small},
%      eje/.style={->},
%      proyeccion/.style={black!20, line width=0.4},
%      longitud/.style={line width=.5pt,\colorV},
%      textooriginal/.style={\colorTorig},
%      pcirculo/.style={fill=\colorPorig, draw=black},
%      angulogirado/.style={fill=\colorAngGreen, draw=\colorGirodos},
%      background/.style={%
%        line width=\bgborderwidth,
%        draw=\bgbordercolor,
%        fill=\bgcolor,
%      },
%      ]
%      % COORDENADAS
%      % Origen
%      \coordinate (O) at (0,0);
%      % Extremo izdo del eje x e inferior del eje y
%      \coordinate (xini) at (\XMLONG cm,0);
%      \coordinate (yini) at (0,\YMLONG cm);
%      % Eje x
%      \path[save path=\ejex] (xini) -- coordinate[pos=0.5] (xtexto) (\XPLONG cm, 0)
%      coordinate (xfin);
%      % Eje y
%      \path[save path=\ejey] (yini) -- coordinate[pos=0.6] (ytexto) (0, \YPLONG cm)
%      coordinate (yfin);
%      % Vector
%      \path[save path=\r] (O) -- coordinate[pos=0.5] (pmid) (\PANG:\PMOD cm)
%      coordinate (P);
%      % Proyección de P en los ejes
%      \path[save path=\proyx] (P) |- (xfin);
%      \path[save path=\proyy] (P) -| (yfin);
%
%      % DIBUJOS
%      % Ángulo theta
%      \path (xfin) -- (O) -- (P) pic
%      [-{Latex[width=2.5pt,length=3.7pt]},angulogirado,
%      "\footnotesize $\theta$",angle radius=20pt, angle eccentricity=0.68]
%      {angle = xfin--O--P};
%      % Proyección de P en los ejes
%      \draw[proyeccion, use path=\proyx];
%      \draw[proyeccion, use path=\proyy];
%      % Ejes
%      % x
%      \draw[eje, use path=\ejex];
%      \node[right] at (xfin) {$x$};
%      %\node[below] at (xtexto) {$r\cos\theta$};
%      % y
%      \draw[eje, use path=\ejey];
%      \node[above] at (yfin) {$y$};
%      %\node[rotate=90, left=8pt] at (ytexto) {$r\sin\theta$};
%      % Coordenada r y punto  P
%      \draw[longitud, use path=\r];
%      \node[above left=-3pt and -1pt] at (pmid) {$r$};
%      \filldraw[fill=green, draw=black] (P) circle[radius=1.2pt];
%      \node[above] at (P) {$P$};
%
%      % Origen
%      \filldraw (O) circle [radius=.2pt];
%      % Fondo amarillo
%      \coordinate (SW) at ($(current bounding box.south west) + (-\HORZ cm,-\VERT cm)$);
%      \coordinate (NE) at ($(current bounding box.north east) + (\HORZ cm,\VERT cm)$);
%      \begin{scope}[on background layer]
%        \draw[background] (SW) rectangle (NE);
%      \end{scope}
%    \end{tikzpicture}
%  \caption{Coordenadas polares $r$ y $\theta$ de un punto $P$ del plano.}
%  \label{fig:R2-coord_polares}
%\end{minipage}
%\hspace{1em}
%\begin{minipage}{0.40\linewidth}
%    % Escala
%    \def\scl{1.20}
%    % 
%    \pgfmathsetmacro{\EJESEXTRA}{-0.2}
%    \pgfmathsetmacro{\EJESLONG}{2.3}
%    % Eje x
%    \pgfmathsetmacro{\XMLONG}{\EJESEXTRA}
%    \pgfmathsetmacro{\XPLONG}{\EJESLONG}
%    % Eje y
%    \pgfmathsetmacro{\YMLONG}{\EJESEXTRA}
%    \pgfmathsetmacro{\YPLONG}{\EJESLONG}
%    % Vector P
%    \pgfmathsetmacro{\PMOD}{2.5}
%    \pgfmathsetmacro{\PANG}{40}
%    % Fondo
%    \pgfmathsetmacro{\HORZ}{0.25}
%    \pgfmathsetmacro{\VERT}{0.25}
%    % 
%    \centering
%    \begin{tikzpicture}[%
%      scale=\scl,
%      every node/.style={black,font=\small},
%      eje/.style={->},
%      proyeccion/.style={black!20, line width=0.4},
%      vector/.style={-{Latex[round]}, shorten >=1.2pt, line width=.8pt,\colorV},
%      textooriginal/.style={\colorTorig},
%      pcirculo/.style={fill=\colorPorig, draw=black},
%      angulogirado/.style={fill=\colorAngGreen, draw=\colorGirodos},
%      background/.style={%
%        line width=\bgborderwidth,
%        draw=\bgbordercolor,
%        fill=\bgcolor,
%      },
%      ]
%      % COORDENADAS
%      % Origen
%      \coordinate (O) at (0,0);
%      % Extremo izdo del eje x e inferior del eje y
%      \coordinate (xini) at (\XMLONG cm,0);
%      \coordinate (yini) at (0,\YMLONG cm);
%      % Eje x
%      \path[save path=\ejex] (xini) -- coordinate[pos=0.5] (xtexto) (\XPLONG cm, 0)
%      coordinate (xfin);
%      % Eje y
%      \path[save path=\ejey] (yini) -- coordinate[pos=0.6] (ytexto) (0, \YPLONG cm)
%      coordinate (yfin);
%      % Vector
%      \path[save path=\vector] (O) -- coordinate[pos=0.5] (pmid) (\PANG:\PMOD cm)
%      coordinate (P);
%      % Proyección de P en los ejes
%      \path[save path=\proyx] (P) |- (xfin);
%      \path[save path=\proyy] (P) -| (yfin);
%
%      % DIBUJOS
%      % Ángulo theta
%      \path (xfin) -- (O) -- (P) pic
%      [-{Latex[width=2.5pt,length=3.7pt]},angulogirado,
%      "\footnotesize $\theta$",angle radius=20pt, angle eccentricity=0.68]
%      {angle = xfin--O--P};
%      % Proyección de P en los ejes
%      \draw[proyeccion, use path=\proyx];
%      \draw[proyeccion, use path=\proyy];
%      % Ejes
%      % x
%      \draw[eje, use path=\ejex];
%      \node[right] at (xfin) {$x$};
%      \node[below] at (xtexto) {$r\cos\theta$};
%      % y
%      \draw[eje, use path=\ejey];
%      \node[above] at (yfin) {$y$};
%      \node[rotate=90, left=8pt] at (ytexto) {$r\sin\theta$};
%      % Vector y punto  P
%      \draw[vector, use path=\vector];
%      \node[above left=-3pt and -1pt] at (pmid) {$\vvv{r}$};
%      \filldraw[fill=green, draw=black] (P) circle[radius=1.2pt];
%      \node[above] at (P) {$P$};
%
%      % Origen
%      \filldraw (O) circle [radius=.2pt];
%      % Fondo amarillo
%      \coordinate (SW) at ($(current bounding box.south west) + (-\HORZ cm,-\VERT cm)$);
%      \coordinate (NE) at ($(current bounding box.north east) + (\HORZ cm,\VERT cm)$);
%      \begin{scope}[on background layer]
%        \draw[background] (SW) rectangle (NE);
%      \end{scope}
%    \end{tikzpicture}
%    \caption{Vector de posición $\vvv{r}$, y coordenadas $x$ e $y$ en función de $r$ y
%      $\theta$.}
%  \label{fig:R2:vector-r-R2}
%\end{minipage}
%\end{figure}
%
%\section{Relación con las coordenadas cartesianas}
%De la figura \ref{fig:R2:vector-r-R2} se desprenden las siguientes relaciones
%\begin{subequations}
%  \begin{align}\label{eq:R2-x-rtheta}
%    x &= r\cos\theta\\
%    \label{eq:R2-y-rtheta}
%    y &= r\sin\theta
%  \end{align}
%\end{subequations}
%y sus inversas
%\begin{subequations}
%  \begin{align}\label{eq:R2-r-xy}
%    r &= \sqrt{x^2 + y^2}\\
%    \label{eq:R2-theta-xy}
%    \theta &= \arctan\left(\frac{y}{x}\right)
%  \end{align}
%\end{subequations}
%Nótese que el ángulo $\theta$ no está definido en el origen. Por eso en este sistema el
%origen es un polo (las coordenadas no están completamente definidas).
%
%\section{Bases, vectores unitarios, sus productos escalares y derivadas}
%El vector $\vvv{r}$ que representa cualquier punto de $\symbb{R}^2$ se puede escribir, en
%la base natural de coordenadas cartesianas, como
%\begin{equation}
%  \vvv{r} = r\cos\theta\,\uvec{\i} + r\sin\theta\,\uvec{\j}
%\end{equation}
%\subsubsection{Derivadas parciales de la posición}
%La base natural en polares está formada por los vectores
%$B = \{\vvv{e}_r, \vvv{e}_{\theta}\}$ ---no necesariamente unitarios--- que llevan
%información del cambio de posición al variar cada coordenada forman una base
%$\set{\vvv{e}_r, \vvv{e}_\theta}$ y se calculan derivando parcialmente el vector
%$\vvv{r}$ con respecto de las variables $r$ y $\theta$
%\begin{subequations}
%\begin{align}
%  \vvv{e}_r
%  &=
%    \dfrac{\partial\vvv{r}}{\partial r}
%    = \cos\theta\,\uvec{\i} + \sin\theta\,\uvec{\j}
%    \longrightarrow
%    |\vvv{e}_r| = 1\\
%  \vvv{e}_\theta
%  &=
%    \dfrac{\partial\vvv{r}}{\partial\theta}
%    = -r\sin\theta\,\uvec{\i} + r\cos\theta\,\uvec{\j}
%    \longrightarrow
%    |\vvv{e}_\theta| = r\\    
%\end{align}
%\end{subequations}
%Estos vectores se utilizan en geometría diferencial y son los que se utilizarán para
%calcular la métrica en coordenadas polares. Nótese que, en particular, $\vvv{e}_\theta$ no
%es un vector unitario, su módulo vale $r$.
%
%\subsubsection{Vectores unitarios}
%En física estamos interesados en bases unitarias. La base unitaria en polares la
%representaremos como $B=\set{\xhat{r}, \xhat{\theta}}$.
%Calculamos los vectores unitarios en polares en función de la base cartesiana
%{\small
%\begin{subequations}
%  \begin{align}\label{eq:R2-versor-r}
%    \xhat{r}
%    &= \dfrac{\vvv{e}_r}{|\vvv{e}_r|}
%    =\dfrac{\cos\theta\,\uvec{\i}+\sin\theta\,\uvec{\j}}
%    {\sqrt{\cos^2\theta+\sin^2\theta}}
%    = \dfrac{\cos\theta\,\uvec{\i}+\sin\theta\,\uvec{\j}}{1}
%      = \cos\theta\,\uvec{\i} + \sin\theta\,\uvec{\j}\\
%    \label{eq:R2-versor-theta}
%  \xhat{\theta}
%  &= \dfrac{\vvv{e}_\theta}{|\vvv{e}_\theta|}
%  = \dfrac{-r\sin\theta\,\uvec{\i}+\cos\theta\,\uvec{\j}}
%  {\sqrt{r^2(\cos^2\theta+\sin^2\theta)}}
%  = \dfrac{-r\sin\theta\,\uvec{\i}+r\cos\theta\,\uvec{\j}}{r}
%  = -\sin\theta\,\uvec{\i} + \cos\theta\,\uvec{\j}
%  \end{align}
%\end{subequations}
%}
%Resumimos este cambio de base y su inversa en forma matricial
%\begin{subequations}
%\begin{align}
%  \begin{pmatrix}
%    \xhat{r}\\
%    \xhat{\theta}
%  \end{pmatrix}
%  &=
%  \begin{pmatrix}
%    \cos\theta & \sin\theta\\
%    -\sin\theta & \cos\theta
%  \end{pmatrix}
%  \,
%  \begin{pmatrix}
%    \uvec{\i}\\
%    \uvec{\j}
%  \end{pmatrix}\\
%\begin{pmatrix}
%    \uvec{\i}\\
%    \uvec{\j}
%  \end{pmatrix}
%  &=
%  \begin{pmatrix}
%    \cos\theta & -\sin\theta\\
%    \sin\theta & \cos\theta
%  \end{pmatrix}
%  \,
%  \begin{pmatrix}
%    \xhat{r}\\
%    \xhat{\theta}
%  \end{pmatrix} 
%\end{align}
%\end{subequations}
%donde se ha utilizado la inversa de la matriz de cambio de base, que por ser ortogonal,
%la inversa es su transpuesta.
%
%Nótese que los vectores de la base natural $\vvv{e}_r$ y $\vvv{e}_\theta$ se pueden
%escribir en función de los unitarios como
%\begin{subequations}
%\begin{align}
%  \vvv{e}_r &= \xhat{r}\\
%  \vvv{e}_\theta &= r\xhat{\theta}
%\end{align}
%\end{subequations}
%
%\subsubsection{Productos escalares}
%\begin{itemize}
%\item De los vectores unitarios
%{\small
%\begin{subequations}
%  \begin{align}
%    \xhat{r}\cdot\xhat{r}
%    &=
%      (\cos\theta\,\uvec{\i}+\sin\theta\,\uvec{\j})
%      \cdot
%      (\cos\theta\,\uvec{\i}+\sin\theta\,\uvec{\j})
%      = \cos^2\theta + \sin^2\theta = 1\\
%    \xhat{r}\cdot\xhat{\theta}
%    &=
%      \xhat{\theta}\cdot\xhat{r}
%      = (\cos\theta\,\uvec{\i}+\sin\theta\,\uvec{\j})
%      \cdot
%      (-\sin\theta\,\uvec{\i}+\cos\theta\,\uvec{\j})
%      = -\sin\theta\cos\theta + \sin\theta\cos\theta = 0\\
%    \xhat{\theta}\cdot\xhat{\theta}
%    &=
%      (-\sin\theta\,\uvec{\i}+\cos\theta\,\uvec{\j})
%      \cdot
%      (-\sin\theta\,\uvec{\i}+\cos\theta\,\uvec{\j})
%      = \sin^2\theta + \cos^2\theta = 1
%  \end{align}
%\end{subequations}
%}
%
%\item De los vectores $\vvv{e}_r$ y $\vvv{e}_\theta$
%\begin{subequations}
%  \begin{align}\label{eq:R2-erer-productoescalar}
%    \vvv{e}_r\cdot\vvv{e}_r
%    &=
%      \xhat{r}\cdot\vvv{r} = 1\\
%    \label{eq:R2-eretheta-productoescalar}
%    \vvv{e}_r\cdot\vvv{e}_\theta
%    &=
%      \vvv{e}_\theta\cdot\vvv{e}_r
%      = \xhat{r}\cdot r\xhat{\theta} = r \xhat{r}\cdot\xhat{\theta} = 0\\
%    \label{eq:R2-ethetaetheta-productoescalar}
%    \vvv{e}_r\cdot\vvv{e}_\theta
%    &=
%      \vvv{e}_\theta\cdot\vvv{e}_\theta
%      = r\xhat{\theta}\cdot r\xhat{\theta}
%      = r^2 \xhat{\theta}\cdot\xhat{\theta} = r^2\\    
%  \end{align}
%\end{subequations}
%\end{itemize}
%
%\subsubsection{Derivadas de los vectores unitarios}
%Para terminar, nos interesan las derivadas de los vectores $\xhat{r}$ y $\xhat{\theta}$
%con respecto de $r$
%\begin{subequations}
%\begin{align}\label{eq:R2-partial-versor_r-r}
%  \dfrac{\partial\xhat{r}}{\partial r}
%  &=
%  \dfrac{\partial}{\partial r}(\cos\theta\,\uvec{\i}+\sin\theta\,\uvec{\j})
%  = 0\\
%  \label{eq:R2-partial-versor_theta-r}
%  \dfrac{\partial\xhat{\theta}}{\partial r}
%  &=
%  \dfrac{\partial}{\partial r}(-\sin\theta\,\uvec{\i}+\cos\theta\,\uvec{\j})
%  = 0
%\end{align}
%\end{subequations}
%Derivadas con respecto de $\theta$
%\begin{subequations}
%\begin{align}\label{eq:R2-partial-versor_r-theta}
%  \dfrac{\partial\xhat{r}}{\partial\theta}
%  &=
%  \dfrac{\partial}{\partial r}(\cos\theta\,\uvec{\i}+\sin\theta\,\uvec{\j})  
%  = -\sin\theta\,\uvec{\i} + \cos\theta\,\uvec{\j} = \xhat{\theta}\\
%  \label{eq:R2-partial-versor_theta-theta}
%  \dfrac{\partial\xhat{\theta}}{\partial\theta}
%  &=
%  \dfrac{\partial}{\partial r}(-\sin\theta\,\uvec{\i}+\cos\theta\,\uvec{\j})  
%  = -\cos\theta\,\uvec{\i} - \sin\theta\,\uvec{\j} = -\xhat{r}
%\end{align}
%\end{subequations}
%
%\section{Métrica y factores de escala}
%Vamos a calcular la métrica y los factores de escala de varias formas:
%\begin{enumerate}
%\item Directamente, utilizando los productos escalares
%  \eqref{eq:R2-erer-productoescalar}, \eqref{eq:R2-eretheta-productoescalar} y
%   \eqref{eq:R2-ethetaetheta-productoescalar}
%  \begin{equation}\label{eq:R2-metrica-polares}
%    \mmm{g}_{ij}
%    = \begin{pmatrix}
%      g_{11} & g_{12}\\
%      g_{21} & g_{22}
%      \end{pmatrix}
%    = \begin{pmatrix}
%      \vvv{e}_r\cdot\vvv{e}_r & \vvv{e}_r\cdot\vvv{e}_\theta\\
%      \vvv{e}_\theta\cdot\vvv{e}_r & \vvv{e}_\theta\cdot\vvv{e}_\theta
%    \end{pmatrix}
%    = \begin{pmatrix}
%      1 & 0\\
%      0 & r^2
%    \end{pmatrix}
%  \end{equation}
%  
%\item Empezamos calculando las diferenciales $dx$ y $dy$
%  \begin{subequations}
%    \begin{align}\label{eq:R2-dx}
%      dx
%      &=
%        \dfrac{\partial x}{\partial r} dr + \dfrac{\partial x}{\partial\theta} d\theta
%        = \cos\theta dr - r\sin\theta d\theta\\
%      \label{eq:R2-dy}
%      dy
%      &=
%        \dfrac{\partial y}{\partial r} dr + \dfrac{\partial y}{\partial\theta} d\theta
%        = \sin\theta dr + r\cos\theta d\theta
%    \end{align}
%  \end{subequations}
%
%La métrica se puede deducir a partir del cuadrado de la distancia recorrida en el
%espacio al cambiar las coordenadas
%{\footnotesize
%  \begin{align*}
%    ds^2
%    &=
%      dx^2 + dy^2
%      = (\cos\theta dr - r\sin\theta d\theta)^2
%      + (\sin\theta dr + r\cos\theta d\theta)^2\\
%    &=
%      \cos^2\theta dr^2
%      + \!r^2\sin^2\theta d\theta^2
%      - \!\cancelout{2r\sin\theta \cos\theta dr d\theta}
%      + \!\sin^2\theta dr^2
%      + \!r^2 \cos^2\theta d\theta^2
%      + \!\cancelout{2r\sin\theta \cos\theta dr d\theta}\\
%    &=
%      (\cos^2\theta + \sin^2\theta) dr^2 + r^2(\sin^2\theta + \cos^2\theta) d\theta^2
%      = dr^2 + r^2 d\theta^2
%  \end{align*}
%}
%donde se han tenido en cuenta las ecuaciones~\eqref{eq:R2-dx} y ~\eqref{eq:R2-dy} y que
%los vectores unitarios $\xhat{r}$ y $\xhat{\theta}$ son ortogonales.
%
%\item De forma intuitiva podemos interpretar la métrica. Veamos
%  \begin{itemize}
%  \item Las variables $r$ y $\theta$ son ortogonales, un cambio en una de ellas no afecta
%    a la otra, como se puede comprobar observando las figuras
%    \ref{fig:R2-polares-intuicion-h-r} y \ref{fig:R2-polares-intuicion-h-theta}.
%    De esta manera, vemos que los elementos no diagonales de la métrica valen cero,
%    $g_{12}=g_{21}=0$.
%  \item Por otro lado, en la figura \ref{fig:R2-polares-intuicion-h-r} observamos que
%    cuando modificamos la variable $r$ en una cierta distancia $\Delta r$, estamos
%    desplazándonos en el espacio la misma distancia $\Delta s$. esto significa que
%    el coeficiente de escala $h_r$ vale la unidad. Esto implica que el elemento
%    $g_{11}$ de la métrica vale lo mismo.
%    \begin{subequations}
%      \begin{align}
%        h_r
%        &=
%          \dfrac{\Delta s}{\Delta r} = 1\\[2pt]
%        g_{11}
%        &=
%          h_x^2 = 1
%      \end{align}
%    \end{subequations}
%
%  \item En cuanto al coeficiente de escala $h_\theta$, cuando se modifica $\theta$ en una
%    cierta cantidad $\Delta\theta$, la distancia que se recorre en el espacio es
%    el ángulo por el radio, $\Delta s = r\Delta\theta$. Además, $g_{22} = r^2$. Ver
%    figura \ref{fig:R2-polares-intuicion-h-theta}
%    \begin{subequations}
%      \begin{align}
%        h_\theta
%        &=
%          \dfrac{\Delta s}{\Delta\theta} = \dfrac{r\Delta\theta}{\Delta\theta} = r\\[2pt]
%        g_{22}
%        &=
%          h_\theta^2 = r^2
%      \end{align}
%    \end{subequations}
%    \begin{figure}[ht]
%      \centering
%      \begin{minipage}{0.40\linewidth}
%        % Esta gráfica muestra cómo afecta un cambio en la coordenada r a la distancia
%        % recorrida en el espacio R^2
%        % Escala
%        \def\scl{1.35}
%        % 
%        \pgfmathsetmacro{\EJESEXTRA}{-0.2}
%        \pgfmathsetmacro{\EJESLONG}{2.5}
%        % Eje x
%        \pgfmathsetmacro{\XMLONG}{\EJESEXTRA}
%        \pgfmathsetmacro{\XPLONG}{\EJESLONG}
%        % Eje y
%        \pgfmathsetmacro{\YMLONG}{\EJESEXTRA}
%        \pgfmathsetmacro{\YPLONG}{\EJESLONG}
%        % Punto P original
%        \pgfmathsetmacro{\PMOD}{1.5}
%        \pgfmathsetmacro{\PANG}{30}
%        % Punto P final conservando el mismo ángulo
%        \pgfmathsetmacro{\PFINMOD}{2.4}
%        \pgfmathsetmacro{\PFINANG}{\PANG}      
%        % Fondo
%        \pgfmathsetmacro{\HORZ}{0.25}
%        \pgfmathsetmacro{\VERT}{0.25}
%        % 
%        \centering
%        \begin{tikzpicture}[%
%          scale=\scl,
%          every node/.style={black,font=\small},
%          eje/.style={->},
%          r/.style={line width=.8pt, black},
%          rtxt/.style={black},
%          rfin/.style={line width=1pt, \colorVred},
%          rfintxt/.style={\colorTchg, rotate=\PANG},
%          textooriginal/.style={\colorTorig},
%          pcirculo/.style={black, radius=1pt},
%          angulogirado/.style={fill=\colorAng, draw=\colorGirodos},
%          background/.style={%
%            line width=\bgborderwidth,
%            draw=\bgbordercolor,
%            fill=\bgcolor,
%          },
%          ]
%          % COORDENADAS
%          % Origen
%          \coordinate (O) at (0,0);
%          % Extremo izdo del eje x e inferior del eje y
%          \coordinate (xini) at (\XMLONG cm,0);
%          \coordinate (yini) at (0,\YMLONG cm);
%          % Eje x
%          \path[save path=\ejex] (xini) -- coordinate[pos=0.5] (xtexto)
%          (\XPLONG cm, 0) coordinate (xfin);
%          % Eje y
%          \path[save path=\ejey] (yini) -- coordinate[pos=0.6] (ytexto)
%          (0, \YPLONG cm) coordinate (yfin);
%          % Punto P original
%          \path[save path=\porig] (O) -- coordinate[pos=0.65] (rorigtxt)
%          (\PANG:\PMOD cm) coordinate (P);
%          % Punto P final
%          \path[save path=\pfin] (P) -- coordinate[pos=0.5] (rfintxt)
%          (\PFINANG:\PFINMOD cm) coordinate (P');
%        
%          % DIBUJOS
%          % Ángulo theta
%          \path (xfin) -- (O) -- (P) pic
%          [-{Latex[width=2.5pt,length=3.7pt]},angulogirado,
%          "\scriptsize $\theta$",angle radius=24pt, angle eccentricity=0.72]
%          {angle = xfin--O--P};
%          % Ejes
%          % x
%          \draw[eje, use path=\ejex];
%          \node[right] at (xfin) {$x$};
%          % y
%          \draw[eje, use path=\ejey];
%          \node[above] at (yfin) {$y$};
%          % Coordenada r original
%          \draw[r, use path=\porig];
%          \node[rtxt, above left=-3pt and -1pt] at (rorigtxt) {\footnotesize $r$};
%          % Vector r final
%          \draw[rfin, use path=\pfin];
%          \node[rfintxt, above] at (rfintxt) {\footnotesize $\Delta s = \Delta r$};
%          % Puntos
%          \fill[pcirculo] (P) circle;
%          \fill[pcirculo] (P') circle;
%          \node[rtxt, below right] at (P) {\scriptsize $P$};
%          \node[rtxt, below right] at (P') {\scriptsize $P'$};
%          % Origen
%          \filldraw (O) circle [radius=.4pt];
%          % Fondo amarillo
%          \coordinate (SW)
%          at ($(current bounding box.south west) + (-\HORZ cm,-\VERT cm)$);
%          \coordinate (NE)
%          at ($(current bounding box.north east) + (\HORZ cm,\VERT cm)$);
%          \begin{scope}[on background layer]
%            \draw[background] (SW) rectangle (NE);
%          \end{scope}
%        \end{tikzpicture}
%        \caption{Al cambiar la coordenada $r$, manteniendo constante $\theta$, se recorre
%          una distancia igual en el espacio $\symbb{R}^2$. Por tanto, $h_r=1$ y
%          $g_{11} = 1$.}
%        \label{fig:R2-polares-intuicion-h-r}
%      \end{minipage}
%      \hspace{1em}
%      \begin{minipage}{0.40\linewidth}
%        % Escala
%        \def\scl{1.35}
%        % 
%        \pgfmathsetmacro{\EJESEXTRA}{-0.2}
%        \pgfmathsetmacro{\EJESLONG}{2.5}
%        % Eje x
%        \pgfmathsetmacro{\XMLONG}{\EJESEXTRA}
%        \pgfmathsetmacro{\XPLONG}{\EJESLONG}
%        % Eje y
%        \pgfmathsetmacro{\YMLONG}{\EJESEXTRA}
%        \pgfmathsetmacro{\YPLONG}{\EJESLONG}
%        % Punto P original
%        \pgfmathsetmacro{\PMOD}{2.0}
%        \pgfmathsetmacro{\PANG}{25}
%        % Punto P final conservando el mismo radio
%        \pgfmathsetmacro{\PFINMOD}{\PMOD}
%        \pgfmathsetmacro{\PFINANG}{70}
%        % Texto en arco
%        \pgfmathsetmacro{\PARCANG}{0.5*(\PFINANG) + \PANG)}
%        \pgfmathsetmacro{\PARCMOD}{\PMOD * 1.2}
%        % Fondo
%        \pgfmathsetmacro{\HORZ}{0.25}
%        \pgfmathsetmacro{\VERT}{0.25}
%        % 
%        \centering
%        \begin{tikzpicture}[%
%          scale=\scl,
%          every node/.style={black,font=\small},
%          eje/.style={->},
%          r/.style={line width=.8pt, black},
%          rtxt/.style={black},
%          rfin/.style={line width=1pt, \colorVred},
%          rfintxt/.style={\colorTchg},
%          textooriginal/.style={\colorTorig},
%          pcirculo/.style={black, radius=1pt},
%          angulo/.style={fill=\colorAng, draw=\colorGirodos},
%          angulogirado/.style={fill=\colorAngPink, draw=\colorGiroRed},
%          arco/.style={line width=1pt, draw=\colorGiroRed},
%          background/.style={%
%            line width=\bgborderwidth,
%            draw=\bgbordercolor,
%            fill=\bgcolor,
%          },
%          ]
%
%          % COORDENADAS
%          % Origen
%          \coordinate (O) at (0,0);
%          % Extremo izdo del eje x e inferior del eje y
%          \coordinate (xini) at (\XMLONG cm,0);
%          \coordinate (yini) at (0,\YMLONG cm);
%          % Eje x
%          \path[save path=\ejex] (xini) -- coordinate[pos=0.5] (xtexto)
%          (\XPLONG cm, 0) coordinate (xfin);
%          % Eje y
%          \path[save path=\ejey] (yini) -- coordinate[pos=0.6] (ytexto)
%          (0, \YPLONG cm) coordinate (yfin);
%          % Punto P original
%          \path[save path=\porig] (O) -- coordinate[pos=0.65] (rorigtxt)
%          (\PANG:\PMOD cm) coordinate (P);
%          % Punto P final
%          \path[save path=\pfin] (O) -- coordinate[pos=0.65] (rfintxt)
%          (\PFINANG:\PFINMOD cm) coordinate (P');
%
%          % DIBUJOS
%          % Ángulo theta
%          \path (xfin) -- (O) -- (P) pic
%          [angulo, "\scriptsize $\theta$",angle radius=24pt, angle eccentricity=0.72]
%          {angle = xfin--O--P};
%          % Incremento de theta
%          \path (P) -- (O) -- (P') pic
%          [-{Latex[width=3.0pt,length=4.5pt]}, angulogirado,
%          "\scriptsize $\Delta\theta$",angle radius=24pt, angle eccentricity=0.72]
%          {angle = P--O--P'};
%          % Ejes
%          % x
%          \draw[eje, use path=\ejex];
%          \node[right] at (xfin) {$x$};
%          % y
%          \draw[eje, use path=\ejey];
%          \node[above] at (yfin) {$y$};
%          % Longitud r original
%          \draw[r, use path=\porig];
%          \node[rtxt, below right=-2pt and 0pt] at (rorigtxt) {\footnotesize $r$};
%          % Longitud r final
%          \draw[r, use path=\pfin];
%          \node[rtxt, above left = -2pt and 0pt] at (rfintxt) {\footnotesize $r$};
%          % Arco recorrido
%          \draw[arco] (P) arc (\PANG:\PFINANG:\PMOD)
%          node[rfintxt,above,midway,sloped] {\footnotesize$\Delta s = r\Delta\theta$};
%        
%          % Puntos
%          \fill[pcirculo] (P) circle;
%          \fill[pcirculo] (P') circle;
%          \node[rtxt, below right=-3pt and 0pt] at (P) {\scriptsize $P$};
%          \node[rtxt, above left=0pt and -3pt] at (P') {\scriptsize $P'$};
%          % Origen
%          \filldraw (O) circle [radius=.4pt];
%          % Fondo amarillo
%          \coordinate (SW)
%          at ($(current bounding box.south west) + (-\HORZ cm,-\VERT cm)$);
%          \coordinate (NE)
%          at ($(current bounding box.north east) + (\HORZ cm,\VERT cm)$);
%          \begin{scope}[on background layer]
%            \draw[background] (SW) rectangle (NE);
%          \end{scope}
%
%        \end{tikzpicture}
%        \caption{Al cambiar la coordenada $\theta$, manteniendo $r$ constante, se  recorre
%          una distancia proporcional a $r$ y a $\Delta\theta$. Por tanto, $h_\theta = r$ y
%          $g_{22} = r^2$.}
%        \label{fig:R2-polares-intuicion-h-theta}
%      \end{minipage}
%    \end{figure}
%  \end{itemize}
%  
%\end{enumerate}
%
%\section{Gradiente}
%Cálculos preliminares. Se utilizan las expresiones \eqref{eq:R2-x-rtheta},
%\eqref{eq:R2-y-rtheta}, \eqref{eq:R2-r-xy} y \eqref{eq:R2-theta-xy}
%{\small
%\begin{subequations}
%  \begin{align}
%    \dfrac{\partial r}{\partial x}
%    &=
%      \dfrac{1}{\cancelout{2}\sqrt{x^2+y^2}}\cdot \cancelout{2}x
%      = \dfrac{\cancelout{r}\cos\theta}{\cancelout{r}} = \cos\theta\\
%    \dfrac{\partial r}{\partial y}
%    &=
%      \dfrac{1}{\cancelout{2}\sqrt{x^2+y^2}}\cdot \cancelout{2}y
%      = \dfrac{\cancelout{r}\sin\theta}{\cancelout{r}} = \sin\theta\\
%    \dfrac{\partial\theta}{\partial x}
%    &=
%      \dfrac{1}{1+\dfrac{y^2}{x^2}}\,\left(-\dfrac{y}{x}\right)
%      = -\dfrac{x^{\cancelout{\scriptstyle{2}}}}{x^2+y^2}\,\dfrac{y}{\cancelout{x}}
%      = -\dfrac{y}{x^2+y^2}
%      = -\dfrac{\cancelout{r}\sin\theta}{r^{\cancelout{2}}}
%      = -\dfrac{\sin\theta}{r}\\
%    \dfrac{\partial\theta}{\partial x}
%    &=
%      \dfrac{1}{1+\dfrac{y^2}{x^2}}\,\dfrac{1}{x}
%      = \dfrac{x^{\cancelout{\scriptstyle{2}}}}{x^2+y^2}\,\dfrac{1}{\cancelout{x}}
%      = \dfrac{y}{x^2+y^2}
%      = \dfrac{\cancelout{r}\cos\theta}{r^{\cancelout{\scriptstyle{2}}}}
%      = \dfrac{\cos\theta}{r}
%  \end{align}
%\end{subequations}
%}
%{\small
%  \begin{subequations}
%    \begin{align*}
%      \dfrac{\partial f}{\partial x}\,\uvec{\i}
%      &=
%        \left(%
%        \dfrac{\partial f}{\partial r}\dfrac{\partial r}{\partial x}
%        + \dfrac{\partial f}{\partial\theta}\dfrac{\partial\theta}{\partial x}
%        \right)
%        \left(%
%        \cos\theta\,\xhat{r}-\sin\theta\,\xhat{\theta}
%        \right)\\
%      &=
%        \left[%
%        \dfrac{\partial f}{\partial r}\cos\theta
%        + \dfrac{\partial f}{\partial\theta}\left(-\dfrac{\sin\theta}{r}\right)
%        \right]
%        \left(%
%        \cos\theta\,\xhat{r} - \sin\theta\,\xhat{\theta}
%        \right)\\
%      &=
%        \cos^2\theta\dfrac{\partial f}{\partial r}\,\xhat{r}
%        -\dfrac{\sin\theta\cos\theta}{r}\dfrac{\partial f}{\partial\theta}\,\xhat{r}
%        - \sin\theta\cos\theta\dfrac{\partial f}{\partial r}\,\xhat{\theta}
%        + \dfrac{\cos^2\theta}{r}\dfrac{\partial f}{\partial\theta}\,\xhat{\theta}
%    \end{align*}
%  \end{subequations}
%}
%{\small
%  \begin{subequations}
%    \begin{align*}
%      \dfrac{\partial f}{\partial y}\,\uvec{\j}
%      &=
%        \left(%
%        \dfrac{\partial f}{\partial r}\dfrac{\partial f}{\partial y}
%        + \dfrac{\partial f}{\partial\theta}\dfrac{\partial\theta}{\partial y}
%        \right)
%        \left(%
%        \sin\theta\,\xhat{r}-\cos\theta\,\xhat{\theta}
%        \right)\\
%      &=
%        \left[%
%        \dfrac{\partial f}{\partial r}\sin\theta
%        + \dfrac{\partial f}{\partial\theta}\dfrac{\cos\theta}{r}
%        \right]
%        \left(%
%        \sin\theta\,\xhat{r} - \cos\theta\,\xhat{\theta}
%        \right)\\
%      &=
%        \sin^2\theta\dfrac{\partial f}{\partial r}\,\xhat{r}
%        + \dfrac{\sin\theta\cos\theta}{r}\dfrac{\partial f}{\partial\theta}\,\xhat{r}
%        + \sin\theta\cos\theta\dfrac{\partial f}{\partial r}\,\xhat{\theta}
%        + \dfrac{\cos^2\theta}{r}\dfrac{\partial f}{\partial\theta}\,\xhat{\theta}
%    \end{align*}
%  \end{subequations}
%}
%
%El gradiente en polares se obtiene sustituyendo las dos expresiones anteriores en
%\eqref{eq:R2-gradiente-cartesianas} y sumando. A continuación presentamos el resultado
%{\small
%  \begin{equation}\label{eq:R2-gradiente-polares}
%    \vvv{\nabla}\cdot f
%    =
%      \dfrac{\partial f}{\partial x}\,\uvec{\i}
%      + \dfrac{\partial f}{\partial y}\,\uvec{\j}
%      = \dfrac{\partial f}{\partial r}\,\xhat{r}
%      + \dfrac{1}{r}\dfrac{\partial f}{\partial\theta}\,\xhat{\theta}
%  \end{equation}
%}
%
%\section{Divergencia}
%La divergencia de un campo vectorial en polares la calcularemos de varias formas pero,
%sin que sirva de precedente, la primera que presentaré no es correcta.
%
%Realizamos un cálculo matricial, similar al realizado en coordenadas cartesianas,
%teniendo en cuenta que en este producto escalar se utiliza la métrica en polares
%\eqref{eq:R2-metrica-polares}
%\begin{equation}
%  \vvv{\nabla}\cdot\vvv{A}
%  \neq
%  \cancelout{%
%    \begin{pmatrix}
%      \partial/\partial r & \partial/\partial\theta      
%    \end{pmatrix}
%    \begin{pmatrix}
%      1 & 0\\
%      0 & r^2
%    \end{pmatrix}
%    \begin{pmatrix}
%      A_r\\
%      A_\theta
%    \end{pmatrix}
%  }
%  = \cancelout{%
%    \dfrac{\partial A_r}{\partial r}
%    + \dfrac{\partial r^2A_\theta}{\partial\theta}
%  }
%\end{equation}
%¿Por qué no es correcto este cálculo, si se ha utilizado la métrica correcta en
%polares, de forma similar a lo que se hizo en cartesianas, ver
%\eqref{eq:R2-divergencia-cartesianas}?
%
%% \hspace{1em}
%No lo es porque, aunque hemos utilizado la métrica, esta no es lo único que cambia
%cuando se pasa de cartesianas a polares, esto es, no solo cambian las distancias
%(la métrica), sino también cambian los vectores de la base con la posición en el plano.
%En este cálculo matricial, y en el que se realizó en cartesianas, estaba implícito el
%que los vectores de la base eran inmutables, lo que es correcto en cartesianas pero no
%lo es en polares, dado que estos cambian con la posición en el plano $\symbb{R}^2$.
%
%En resumen, el sistema de coordenadas polares es curvilíneo ver figura
%\ref{fig:R2-coordenadas-polares}, y esa curvatura hay quetenerla en cuenta.
%
%\begin{figure}[ht]
%  % Escala
%  \def\scl{0.9}
%  % 
%  % Número de círcunferencias
%  \pgfmathsetmacro{\NUMCIRC}{5}
%  % Longitud de un versor
%  \pgfmathsetmacro{\VERSOR}{0.5}
%  % Longitud de línea radial
%  \pgfmathsetmacro{\LINEARADIAL}{\NUMCIRC * \VERSOR + 0.25}
%  % Fondo
%  \pgfmathsetmacro{\HORZ}{0.25}
%  \pgfmathsetmacro{\VERT}{0.25}
%  % 
%  \centering
%  \begin{tikzpicture}[%
%    scale=\scl,
%    every node/.style={black,font=\small},
%    vtexto/.style={font=\footnotesize, pos=1.4},
%    vtheta/.style={-{Latex[round, width=4pt, length=5pt]},
%      line width=0.6pt, \colorVred},
%    vr/.style={vtheta, \colorV},
%    polares/.style={draw=black!25, line width=0.4pt},
%    punto/.style={fill=black, radius=1.25pt},
%    background/.style={%
%      line width=\bgborderwidth,
%      draw=\bgbordercolor,
%      fill=\bgcolor
%    },
%    ]
%    % COORDENADAS
%    % Origen
%    \coordinate (O) at (0,0);
%    
%    % DIBUJOS
%    % Coordenadas polares (Radios de los ángulos)
%    \foreach \angulo / \label
%    in {%
%      0/{0}, 30/{\pi/6}, 60/{\pi/3},
%      90/{\pi/2}, 120/{2\pi/3}, 150/{5\pi/6},
%      180/{\pi}, 210/{7\pi/6}, 240/{4\pi/3},
%      270/{3\pi/2}, 300/{5\pi/3}, 330/{11\pi/6}
%    }
%    {%
%      \draw[polares] (O) -- ++(\angulo:\LINEARADIAL)
%      node[pos=1.12] {\scriptsize $\label$};
%    };
%    \foreach \radio [evaluate=\radio] in {\VERSOR*1,\VERSOR*...,\VERSOR*\NUMCIRC}
%    {%
%      % Circunferencias
%      \draw[polares] (O) circle[radius=\radio cm];
%    };
%    
%    % Vectores unitarios en r=1.0
%    \foreach \radio [evaluate=\radio] in{\VERSOR*1}
%    \foreach \angulo in {30}
%    {%
%      % Vectores unitario radiales
%      \draw[vr] (O) ++(\angulo:\radio) -- ++(\angulo:\VERSOR)
%      node[vtexto, rotate=\angulo-90] {$\xhat{r}$};
%      % Vectores unitarios angulares
%      \draw[vtheta] (O) ++(\angulo:\radio) -- ++(90+\angulo:\VERSOR)
%      node[vtexto, rotate=\angulo-90, \colorTred] {$\xhat{\theta}$};
%      % Posición
%      \fill[punto] (O) ++(\angulo:\radio) circle;
%    };
%    
%    % Vectores unitarios en r=2.0
%    \foreach \radio [evaluate=\radio] in{\VERSOR*2}
%    \foreach \angulo in {150}
%    {%
%      % Vectores unitario radiales
%      \draw[vr] (O) ++(\angulo:\radio) -- ++(\angulo:\VERSOR)
%      node[vtexto, rotate=\angulo-90] {$\xhat{r}$};
%      % Vectores unitarios angulares
%      \draw[vtheta] (O) ++(\angulo:\radio) -- ++(90+\angulo:\VERSOR)
%      node[vtexto, rotate=\angulo-90, \colorTred] {$\xhat{\theta}$};
%      % Posición
%      \fill[punto] (O) ++(\angulo:\radio) circle;
%    };
%    
%    % Vectores unitarios en r=3.0
%    \foreach \radio [evaluate=\radio] in{\VERSOR*3}
%    \foreach \angulo in {210, 330}
%    {%
%      % Vectores unitario radiales
%      \draw[vr] (O) ++(\angulo:\radio) -- ++(\angulo:\VERSOR)
%      node[vtexto, rotate=\angulo-90] {$\xhat{r}$};
%      % Vectores unitarios angulares
%      \draw[vtheta] (O) ++(\angulo:\radio) -- ++(90+\angulo:\VERSOR)
%      node[vtexto, rotate=\angulo-90, \colorTred] {$\xhat{\theta}$};
%      % Posición
%      \fill[punto] (O) ++(\angulo:\radio) circle;
%    };
%    
%    % Vectores unitarios en r=4.0
%    \foreach \radio [evaluate=\radio] in{\VERSOR*4}
%    \foreach \angulo in {30,90,270}
%    {%
%      % Vectores unitario radiales
%      \draw[vr] (O) ++(\angulo:\radio) -- ++(\angulo:\VERSOR)
%      node[vtexto, rotate=\angulo-90] {$\xhat{r}$};
%      % Vectores unitarios angulares
%      \draw[vtheta] (O) ++(\angulo:\radio) -- ++(90+\angulo:\VERSOR)
%      node[vtexto, rotate=\angulo-90, \colorTred] {$\xhat{\theta}$};
%      % Posición
%      \fill[punto] (O) ++(\angulo:\radio) circle;
%    };      
%    
%    % Origen
%    \filldraw (O) circle [radius=.2pt];
%    
%    % Fondo amarillo
%    \coordinate (SW) at ($(current bounding box.south west) + (-\HORZ cm,-\VERT cm)$);
%    \coordinate (NE) at ($(current bounding box.north east) + (\HORZ cm,\VERT cm)$);
%    \begin{scope}[on background layer]
%      \draw[background] (SW) rectangle (NE);
%    \end{scope}
%  \end{tikzpicture}
%  \caption{Coordenadas polares (curvilíneas), donde $r$ representa la distancia al
%    origen (polo) y $\theta$ es el ángulo que forma el vector de posición con el eje
%    $x$, aquí expresado en radianes. Además se muestran algunos vectores unidad
%    $\xhat{r}$ y $\xhat{\theta}$ en distintas posiciones del plano $\symbb{R}^2$.}
%  \label{fig:R2-coordenadas-polares}
%\end{figure}
%¿Entonces, no se puede plantear un enfoque matricial para obtener la divergencia en
%polares? Sí es posible, y lo veremos al final.
%
%\begin{enumerate}
%\item Enfoque basado en la variación de los vectores unitarios.
%  
%  Cuando el operador nabla actúa sobre un campo vectorial como
%  $\vvv{A} = A_r\xhat{r}+A_\theta\xhat{\theta}$, no solo actúa sobre las componentes
%  $A_r$ y $A_\theta$, sino también sobre los vectores unitarios $\xhat{r}$ y
%  $\xhat{\theta}$,
%  ver \eqref{eq:R2-versor-r} y \eqref{eq:R2-versor-theta}. Necesitaremos las derivadas
%  de los versores de la base que se calcularon en \eqref{eq:R2-partial-versor_r-r},
%  \eqref{eq:R2-partial-versor_theta-r}, \eqref{eq:R2-partial-versor_r-theta} y
%  \eqref{eq:R2-partial-versor_theta-theta}, que resumimos a continuación:
%  \begin{itemize}
%  \item Al desplazarnos radialmente, los versores $\xhat{r}$ y $\xhat{\theta}$ no
%    cambian, ver ángulo $\pi/6$\,\unit{\radian} en la figura
%    \ref{fig:R2-coordenadas-polares}
%    \begin{equation}
%      \dfrac{\partial\xhat{r}}{\partial r} = \dfrac{\partial\xhat{\theta}}{\partial r}
%      = 0
%    \end{equation}
%  \item Al rotar infinitesimalmente en sentido antihorario, $\xhat{r}$ cambia según
%    $\xhat{\theta}$, ver por ejemplo, ángulos $\pi/6$\,\unit{\radian} y
%    $\pi/2$\,\unit{\radian} en la figura \ref{fig:R2-coordenadas-polares}
%    \begin{equation}
%      \dfrac{\partial\xhat{r}}{\partial\theta} = \xhat{\theta}
%    \end{equation}
%  \item Al desplazarnos infinitesimalmente en sentido antihorario, $\xhat{\theta}$ cambia
%    según el sentido contrario a $\xhat{r}$, ver por ejemplo, ángulos
%    $\pi/6$\,\unit{\radian} y $\pi/2$\,\unit{\radian} en la
%    figura~\ref{fig:R2-coordenadas-polares}
%    \begin{equation}
%      \dfrac{\partial\xhat{\theta}}{\partial\theta} = -\xhat{r}
%    \end{equation}
%    
%  \item También necesitamos el operador nabla en polares, que podemos obtenerlo a
%    partir del gradiente de un campo escalar en polares \eqref{eq:R2-gradiente-polares}
%    \begin{equation}
%      \vvv{\nabla}
%      = \xhat{r}\,\dfrac{\partial}{\partial r}
%      + \xhat{\theta}\,\dfrac{1}{r}\dfrac{\partial}{\partial\theta}
%    \end{equation}
%  \end{itemize}
%  
%  Con todo lo anterior, planteamos el cálculo de la divergencia de un campo vectorial en
%  polares
%  \[
%    \vvv{\nabla}\cdot\vvv{A}
%    =
%    \left(%
%      \xhat{r}\dfrac{\partial}{\partial r}
%      + \xhat{\theta}\dfrac{\partial}{\partial\theta}
%    \right)
%    \cdot
%    (A_r\xhat{r}+A_\theta\xhat{\theta})
%  \]
%  Obtenemos dos sumandos
%  \begin{equation}\label{eq:R2-divergencia-sumandos}
%    \vvv{\nabla}\cdot\vvv{A}
%    = \xhat{r}\cdot
%    \dfrac{\partial}{\partial r}\left(A_r\xhat{r}+A_\theta\xhat{\theta}\right)
%    + \dfrac{1}{r}\xhat{\theta}\cdot\dfrac{\partial}{\partial\theta}
%    \left(A_r\xhat{r}+A_\theta\xhat{\theta}\right)
%  \end{equation}
%  
%  Desarrollamos el primer sumando, teniendo en cuenta que se anulan, tanto las derivadas
%  parciales con respecto a $r$, como el término que contiene a $\xhat{\theta}$, dado que
%  al multiplicarlo por $\xhat{r}$ dará cero, ya que son ortogonales
%  \begin{align*}
%    \xhat{r}\cdot
%    \dfrac{\partial}{\partial r}\left(A_r\xhat{r}+A_\theta\xhat{\theta}\right)
%    &=
%      \xhat{r}\cdot
%      \left(%
%      \dfrac{\partial A_r}{\partial r}\xhat{r}
%      + \cancelout{A_r\dfrac{\partial\xhat{r}}{\partial r}}
%      + \cancelout{\dfrac{\partial A_\theta}{\partial r}\xhat{\theta}}
%      + \cancelout{A_\theta\dfrac{\partial\xhat{\theta}}{\partial r}}
%      \right)\\
%    &=
%      \dfrac{\partial A_r}{\partial r}\,\xhat{r}\cdot\xhat{r}
%      = \dfrac{\partial A_r}{\partial r}
%  \end{align*}
%  Ahora desarrollamos el último sumando, sustituyendo el valor de las derivadas parciales
%  y teniendo en cuenta que $\xhat{r}$ y $\xhat{\theta}$ son ortogonales
%  \begin{align*}
%    \dfrac{1}{r}\xhat{\theta}\cdot\dfrac{\partial}{\partial\theta}
%    \left(A_r\xhat{r}+A_\theta\xhat{\theta}\right)
%    &=
%      \dfrac{1}{r}\xhat{\theta}\cdot
%      \left(%
%      \cancelout{\dfrac{\partial A_r}{\partial\theta}\xhat{r}}
%      + A_r\dfrac{\partial\xhat{r}}{\partial\theta}
%      + \dfrac{\partial A_\theta}{\partial\theta}\xhat{\theta}
%      + A_\theta\dfrac{\partial\xhat{\theta}}{\partial\theta}
%      \right)\\
%    &= 
%      \dfrac{1}{r}\xhat{\theta}\cdot
%      \left(%
%      A_r\xhat{\theta}
%      + \dfrac{\partial A_\theta}{\partial\theta}\xhat{\theta}
%      + \cancelout{A_\theta\,(-\xhat{r})}
%      \right)\\
%    &=
%      \dfrac{1}{r}
%      \left(%
%      A_r + \dfrac{\partial A_\theta}{\partial\theta}
%      \right)
%      \xhat{\theta}\cdot\xhat{\theta}
%      = \dfrac{1}{r}\dfrac{\partial A_\theta}{\partial\theta}
%      + \dfrac{1}{r}A_r
%  \end{align*}
%  
%  Sustituimos estos resultados en los dos sumandos de la divergencia
%  \eqref{eq:R2-divergencia-sumandos}
%  \[
%    \vvv{\nabla}\cdot\vvv{A}
%    = \dfrac{\partial A_r}{\partial r}
%    + \dfrac{1}{r} A_r
%    + \dfrac{1}{r} \dfrac{\partial A_\theta}{\partial\theta}
%  \]
%  
%  Agrupando términos obtenemos la divergencia en polares
%  \begin{equation}\label{eq:R2-divergencia-polares}
%    \vvv{\nabla}\cdot\vvv{A}
%    = \dfrac{1}{r}\dfrac{\partial}{\partial r}(rA_r)
%    + \dfrac{1}{r}\dfrac{\partial A_\theta}{\partial\theta}
%  \end{equation}
%
%\item Cálculo utilizando la métrica y válido en polares y en espacios curvos
%  \begin{equation}
%    \vvv{\nabla}\cdot\vvv{A}
%    = \dfrac{1}{\sqrt{g}}\,\partial_i\big(\sqrt{g} A^i\big)
%  \end{equation}
%  Aplicando la fórmula, sabiendo que la raíz cuadrada del valor absoluto de la métrica
%  es $g=r$ y que $r$ es independiente de $\theta$
%  \begin{align*}
%    \vvv{\nabla}\cdot\vvv{A}
%    &=
%      \dfrac{1}{r}\,\left[%
%      \partial_r\!\left(rA^r\right) + \partial_\theta\!\left(rA^\theta\right)
%      \right]
%    =
%      \dfrac{1}{r}\,\left[%
%      \partial_r\!\left(rA^r\right) + r\,\partial_\theta A^\theta
%      \right]      
%  \end{align*}
%  donde $A^r$ y $A^\theta$ son las componentes contravariantes del vector. En física
%  normalmente se usan las componentes covariantes $A_r$ y $A_\theta$.
%
%  Calculamos la inversa del tensor métrico, que como es diagonal, la inversa de cada
%  elemento es el recíproco de los elementos diagonales
%  \begin{equation}
%    g_{ij} = \begin{pmatrix}1 & 0\\ 0 & 1/r^2 \end{pmatrix}
%  \end{equation}
%  Para bajar el índice, multiplicamos por la inversa del  tensor métrico $A^i = \mmm{g}^{ij}
%  A_j$.
%  Así, desarrollando el sumatorio
%  \begin{align*}
%    A^r
%    &=
%      g^{rr}A_r + g^{r\theta}A_\theta = 1\cdot A_r + 0\cdot A_\theta = A_r\\
%    A^\theta
%    &=
%      g^{\theta r}A_r + g^{\theta\theta}A_\theta = 0\cdot A_r + \dfrac{1}{r^2} A_\theta
%      = \dfrac{1}{r^2}A_\theta
%  \end{align*}
%  Sustituimos en
%  \begin{align*}
%    \vvv{\nabla}\cdot\vvv{A}
%    = \dfrac{1}{r}\left[%
%    \dfrac{\partial}{\partial r}\!\left(r\,A_r\right)
%    + r\, \partial_\theta\left(\dfrac{1}{r^2}\, A_\theta\right)
%    \right]
%    = \dfrac{1}{r}\left[%
%    \dfrac{\partial}{\partial r}\!\left(r\,A_r\right)
%    + \dfrac{1}{r}\, \partial_\theta A_\theta
%    \right]
%  \end{align*}
%  
%\end{enumerate}
%
%
%\section{Laplaciana} 
%
%\subsection{Elemento de área en coordenadas polares planas}
%\begin{align*}
%  d^2a
%  &=
%    \dfrac{\partial(x,y)}{\partial(r,\theta)} dr\,d\theta
%  = \begin{vmatrix}
%    \partial x/\partial r & \partial x/\partial\theta\\
%    \partial y/\partial r & \partial y/\partial\theta
%  \end{vmatrix}
%  dr\,d\theta
%  = \begin{vmatrix}
%    \cos\theta & -r\sin\theta\\
%    \sin\theta & r\cos\theta
%  \end{vmatrix}
%    dr\,d\theta\\
%  &=
%    (r\cos^2\theta + r\sin^2\theta)\,dr\,d\theta
%    = r (\cos^2\theta + \sin^2\theta)\,dr\,d\theta
%\end{align*}
%donde hemos utilizado las relaciones \eqref{eq:R2-x-rtheta} y \eqref{eq:R2-y-rtheta}.
%Simplificando finalmente, queda
%\begin{equation}
%  d^2a = rdr\,d\theta
%\end{equation}



%%% Local Variables:
%%% mode: latex
%%% TeX-engine: luatex
%%% TeX-master: "../retazosmatematicas.tex"
%%% End:
